%  Informe de Practica Profesional II
%  Compilar con: pdflatex (dos veces para el indice)
%==============================================================
\documentclass[12pt,letterpaper,oneside]{report}

%---------------------- Codificacion e idioma -------------------
\usepackage[utf8]{inputenc}
\usepackage[T1]{fontenc}
\usepackage[spanish,es-noshorthands,es-tabla]{babel}

%---------------------- Fuentes (Times) -------------------------
\usepackage{newtxtext}
\usepackage{amsmath,amssymb}
\usepackage{newtxmath}

%---------------------- Pagina y margenes -----------------------
\usepackage[letterpaper,top=2.5cm,bottom=2.5cm,left=3cm,right=2.5cm]{geometry}
\usepackage{setspace}
%\onehalfspacing                 % interlineado 1.5
\doublespacing    
%\usepackage{indentfirst}        % sangria en primer parrafo
%\setlength{\parindent}{1.25cm}
%\setlength{\parskip}{0pt}

\setlength{\parindent}{0pt}
\setlength{\parskip}{6pt}

%---------------------- Graficos y figuras ----------------------
\usepackage{graphicx}
\usepackage{float}
\graphicspath{{Figuras/}}

%---------------------- Tablas ----------------------------------
\usepackage{booktabs}
\usepackage{array}
\usepackage{multirow}
\usepackage{colortbl}
\usepackage{tabularx}



%---------------------- Encabezados -----------------------------
\usepackage{fancyhdr}
\pagestyle{fancy}
\fancyhf{}
\fancyhead[R]{\thepage}
\renewcommand{\headrulewidth}{0pt}

%---------------------- Titulos de capitulos/secciones ----------
\usepackage{titlesec}
\titleformat{\chapter}[block]
  {\normalfont\Large\bfseries}{Capítulo \thechapter:}{0.6em}{}
\titlespacing*{\chapter}{0pt}{-20pt}{20pt}
\titleformat{\section}{\normalfont\large\bfseries}{\thesection}{0.8em}{}
\titleformat{\subsection}{\normalfont\normalsize\bfseries}{\thesubsection}{0.8em}{}

%---------------------- Captions --------------------------------
\usepackage{caption}
\captionsetup[table]{position=above,skip=6pt}
\captionsetup[figure]{position=below,skip=10pt}
\captionsetup{font=small,labelfont=bf}

%---------------------- Comando util ----------------------------
% Inserta la figura si existe el archivo; si no, deja un recuadro
% indicativo con el nombre del fichero que falta.
\newcommand{\figura}[4]{%
  \begin{figure}[H]
    \centering
    \IfFileExists{Figuras/#1}{%
      \includegraphics[width=#2]{#1}%
    }{%
      \fbox{\parbox[c][6cm][c]{0.8\textwidth}{\centering
        \textit{[Insertar aquí la figura: \texttt{ \detokenize{Figuras/#1}}]}}}}%
    \caption{#3}\label{#4}
  \end{figure}%
}

%==============================================================
%                         BIBLIOGRAFIA
%==============================================================
%\usepackage{csquotes}
\usepackage[backend=biber, style=ieee]{biblatex}
\addbibresource{referencias.bib}
%---------------------- Hipervinculos ---------------------------
\usepackage[plainpages=false,pdfpagelabels ,hidelinks]{hyperref}

\hyphenpenalty=10000
\exhyphenpenalty=10000
%\usepackage{microtype}
\widowpenalty=10000
\clubpenalty=10000
%==============================================================
\usepackage{pdflscape}
\usepackage{rotating}
%==============================================================
\begin{document}
\pagenumbering{roman}


%==============================================================
%                         PORTADA
%==============================================================
\begin{titlepage}
\centering
{\large \textbf{Universidad Tecnológica de La Habana}}\\[2pt]
{\large \textbf{``José Antonio Echeverría''}}\\[2pt]
{\large \textbf{Facultad de Ingeniería Informática}}\\[40pt]

\IfFileExists{Figuras/isotipo.png}{%
  \includegraphics[width=3.2cm]{isotipo.png}\\[40pt]%
}{%
  \fbox{\parbox[c][3.2cm][c]{3.2cm}{\centering\textit{[Logo: Figuras/isotipo.png]}}}\\[40pt]%
}

{\Large \textbf{Informe de Práctica Profesional II}}\\[30pt]

{\large \textbf{Implementación de nuevo algoritmo de agrupamiento basado en densidad en la Biblioteca de Heurísticas de Asignación para Problemas de Planificación de Rutas de Vehículos con Múltiples Depósitos}}\\[45pt]

\begin{flushleft}
\hspace{2cm}{\large \textbf{Autor:}}\\[6pt]
\hspace{2cm}{\large Maykel Kevin Díaz Menéndez}\\[25pt]
\hspace{2cm}{\large \textbf{Tutores:}}\\[6pt]
\hspace{2cm}{\large Ing.\ Eric Ramos Aragón}\\
\hspace{2cm}{\large Ing.\ Ananda de la Caridad Morales Morales}
\end{flushleft}

\vfill
{\large \textbf{La Habana, Cuba}}\\[6pt]
{\large \textbf{Septiembre 2026}}
\end{titlepage}

%==============================================================
%                         RESUMEN
%==============================================================
\chapter*{Resumen}

El Problema de Planificación de Rutas de Vehículos con Múltiples Depósitos (MDVRP) constituye una variante compleja dentro de los problemas de optimización combinatoria. En la práctica, este problema suele resolverse dividiéndolo en dos etapas: asignar los clientes a los depósitos y construir las rutas desde cada depósito. La Biblioteca de Heurísticas de Asignación para Problemas de Planificación de Rutas de Vehículos con Múltiples Depósitos (BHAVRP) es un componente de software dise\~nado para resolver esta primera fase de asignación. La biblioteca presenta una arquitectura modular que encapsula algoritmos de agrupamiento clasificados en particionales, jerárquicos y basados en densidad, incluye un sistema de tratamiento de excepciones y un conjunto de métricas de evaluación internas. Actualmente, BHAVRP  cuenta con un solo algoritmo de agrupamiento basado en densidad. Se propone la incorporación de un nuevo algoritmo de agrupamiento basado en densidad HDBSCAN. Este algoritmo de agrupamiento es adaptado para resolver la fase de asignación de clientes a depósitos incorporando restricciones de capacidad de los depósitos y mecanismos de asignación controlada usando como base la implemetación de la biblioteca JSAT. Por último, se realiza un estudio experimental que permite garantizar el correcto funcionamiento de los módulos implementados, así como realizar un análisis de correlación entre las metricas internas y las rutas obtenidas por los clústeres. Para ello, se utilizan los criterios definidos por las métricas de evaluación que ya tiene BHAVRP y la Biblioteca de Heurísticas de Construcción para Problemas de Planificación de Rutas de Vehículos (BHCVRP).

\vspace{12pt}
\noindent\textbf{Palabras clave:} algoritmos de agrupamiento, BHAVRP, HDBSCAN, MDVRP.

%==============================================================
%                         ABSTRACT
%==============================================================
\chapter*{Abstract}

The Multi-Depot Vehicle Routing Problem (MDVRP) constitutes a complex variant within combinatorial optimization problems. In practice, this problem is often solved by dividing it into two stages: assigning customers to depots and building routes from each depot. The Assignment Heuristics Library for Multi-Depot Vehicle Routing Problems (BHAVRP) is a software component designed to solve this first assignment phase. The library features a modular architecture that encapsulates clustering algorithms classified into partitional, hierarchical, and density-based, includes an exception handling system, and a set of internal evaluation metrics. Currently, BHAVRP has only one density-based clustering algorithm. The incorporation of a new density-based clustering algorithm, HDBSCAN, is proposed. This clustering algorithm is adapted to solve the customer-to-depot assignment phase by incorporating depot capacity constraints and controlled assignment mechanisms, using the implementation from the JSAT library as a basis. Finally, an experimental study is conducted to ensure the correct functioning of the implemented modules, as well as to perform a correlation analysis between the internal metrics and the routes obtained by the clusters. For this purpose, the criteria defined by the evaluation metrics already present in BHAVRP and the Construction Heuristics Library for Vehicle Routing Problems (BHCVRP) are used.

\vspace{12pt}
\noindent\textbf{Keywords:} BHAVRP, clustering algorithms, HDBSCAN, MDVRP.

%==============================================================
%                       INDICES
%==============================================================
\tableofcontents
\listoffigures
\listoftables

\clearpage
\pagenumbering{arabic}
\setcounter{page}{1}

%==============================================================
%                      INTRODUCCION
%==============================================================
\chapter*{Introducción}
\addcontentsline{toc}{chapter}{Introducción}
\markboth{Introducción}{}

La gestión de la distribución, el transporte y la planificación de rutas son temas de gran importancia que han sido abordados a nivel internacional. Son ramas de la logística donde yacen numerosas oportunidades de mejora y por esta razón ha recibido especial atención motivando un intenso trabajo teórico en el desarrollo de modelos y algoritmos. En este contexto surgen los Problemas de Planificación de Rutas de Vehículos (\textit{Vehicle Routing Problems}, VRP) \parencite{colomasalazarGestionDistribucionPlanificacion2022}.

En \parencite{g.b.dantzigTruckDispatchingProblem1959} se introduce el \textit{``Truck Dispatching Problem''}, modelando cómo una flota de camiones homogéneos podría satisfacer la demanda de petróleo de varias estaciones de servicio desde un depósito central y con una distancia mínima recorrida. Cinco a\~nos después, en \parencite{clarkeSchedulingVehiclesCentral1964} se generaliza este a un problema de optimización lineal que se encuentra comúnmente en el ámbito de la logística y el transporte. El VRP es un problema NP-duro\parencite{aroraComputationalComplexityModern2009} por lo que se consideran muy difíciles de resolver. Los problemas de la vida real son considerablemente más grandes y complejos que los escenarios académicos, por lo que las heurísticas y metaheurísticas son necesarias para obtener buenos resultados. Estos algoritmos no garantizan encontrar la solución óptima, pero sí proporcionan tiempos computacionales razonables. La cantidad de métodos de solución introducidos en la literatura académica ha crecido en las últimas décadas  \parencite{braekersVehicleRoutingProblem2016}.

El Problema de Planificación de Rutas de Vehículos con Múltiples Depósitos (\textit{Multi-Depot Vehicle Routing Problems}, MDVRP) es una variante del VRP clásico que incorpora varios depósitos con una ubicación predefinida. Cada depósito cuenta con una flota de vehículos de capacidad homogénea que tiene como objetivo satisfacer la demanda de un conjunto de clientes dispersos geográficamente. Cada vehículo debe comenzar y terminar su ruta en el mismo depósito, y cada cliente debe ser visitado una única vez. Teniendo en cuenta la distancia entre la ubicación de los clientes y depósitos, el MDVRP busca asignar un cliente a cada depósito y se construye una ruta para cada vehículo dentro de un depósito para atender dichos clientes. El objetivo al crear la ruta es reducir al mínimo el costo total del recorrido \parencite{mateosDESARROLLOMETODOHIBRIDO2013}.

Durante un largo período de tiempo, los métodos heurísticos basados en la búsqueda tabú han sido los más empleados en la literatura \parencite{mateosDESARROLLOMETODOHIBRIDO2013}. Sin embargo, este problema genera una gran cantidad de combinaciones que se refleja en los altos tiempos computacionales para obtener la solución. El potencial de las técnicas de análisis de clústeres para la solución de los MDVRP ha sido reconocido por autores como \parencite{g.b.dantzigTruckDispatchingProblem1959}, \parencite{bermudezAplicacionTecnicasClusterizacion}, \parencite{barretoUsingClusteringAnalysis2007} y \parencite{mehrjerdiUsingGreedyClustering2013}. El análisis de clústeres estudia la división de objetos en grupos basado en una o varias características. Un grupo se define como una zona con muchos objetos parecidos muy cercanos entre sí, que está separada de otros grupos por zonas donde casi no hay objetos. Esta definición de grupo es bastante acorde a la agrupación de clientes por lo que se propone usar el análisis de clústeres en la solución del MDVRP, específicamente en la fase de asignación de clientes a depósitos \parencite{bermudezAplicacionTecnicasClusterizacion}. En \parencite{jainAlgorithmsClusteringData1988} se describe la técnica ''divide y vencerás'', utilizada para reducir la complejidad computacional en el empleo heurísticas para resolver el MDVRP.

Sin embargo, aún existen desafíos en cuanto a la eficiencia computacional debido a los altos tiempos y la calidad de las agrupaciones dependiendo del escenario \parencite{bermudezAplicacionTecnicasClusterizacion}. En la literatura se destacan algoritmos como \textit{Density-Based Spatial Clustering of Applications with Noise} (DBSCAN)\parencite{khanDBSCANPresentFuture2014} y \textit{K-Means}\parencite{macqueen1967some} que han sido probados y utilizados en el ámbito del MDVRP, mostrando su utilidad para la fase de asignación de clientes a depósitos  \parencite{padulohRobustOptimizingReverse2023, luo2014multi, geetha2013nested}.

En este contexto, la Universidad Tecnológica de la Habana ``José Antonio Echeverría'' (CUJAE) ha desarrollado la Biblioteca de Heurísticas de Asignación para Problemas de Planificación de Rutas de Vehículos con Múltiples Depósitos (BHAVRP). Este componente está dise\~nado para resolver la fase de asignación en MDVRP utilizando una arquitectura modular y 20 algoritmos repartidos en heurísticas clásicas de asignación y algoritmos de agrupamiento de Minería de Datos adaptados al contexto de los MDVRP \parencite{GitHubVrpsolutionsBHAVRP2026}. Esta alternativa gratuita y de código abierto supone una ventaja frente a otras herramientas que no tienen una colección de algoritmos tan extensa o son software comercial.

Esta biblioteca, en la categoría de algoritmos de agrupamiento, incluye algoritmos particionales, algoritmos jerárquicos y un algoritmo basado en densidad. En una investigación anterior \parencite{ericramosaragonIncorporacionModuloMetricas2025}, se presentó los resultados de algoritmos de agrupamiento particionales y heurísticas clásicas de asignación. Sin embargo, estos algoritmos presentaron una tendencia a asignar clientes a depósitos sin considerar adecuadamente la densidad geográfica, lo que generó clústeres con clientes muy alejados del depósito asignado. Esto se conoce en la literatura como ruido o asignaciones atípicas, lo cual puede incrementar los costos de transporte y afecta la calidad de las rutas posteriores. La literatura reconoce que los algoritmos basados en densidad, como DBSCAN, son capaces de identificar y aislar estos puntos atípicos \parencite{khanDBSCANPresentFuture2014}, ya que este tipo de algoritmos parten de la idea de que un clúster es una región de alta densidad de puntos separadas de otras regiones densas por zonas de más baja densidad.

Debido a la situación problemática antes descrita, se identifica como problema a resolver que los algoritmos de agrupamiento particionales y las heurísticas clásicas de asignación implementadas en BHAVRP generan clústeres con presencia de ruido, lo que impacta negativamente en la calidad de las rutas posteriores.

En este contexto, el objeto de estudio de este trabajo es la fase de asignación dentro de los MDVRP y la evaluación de la calidad de los algoritmos de agrupamiento. El campo de acción se enfoca en la incorporación y adaptación de algoritmos de agrupamiento basados en densidad en la fase de asignación de clientes a depósitos de BHAVRP.

Para dar solución al problema planteado se traza como objetivo general: Desarrollar una nueva versión de BHAVRP que incluya más algoritmos de agrupamiento basados en densidad.
Para dar cumplimiento al objetivo general se definieron los siguientes objetivos específicos con sus respectivas tareas:

\begin{enumerate}
  \item Identificar algoritmos de agrupamiento utilizados para resolver la fase de asignación de los problemas de planificación de rutas de vehículos con múltiples depósitos.
  \begin{itemize}
    \item Realizar búsqueda bibliográfica de las características de los problemas de planificación de rutas de vehículos con múltiples depósitos y sus aplicaciones prácticas.
    \item Realizar búsqueda bibliográfica de las propuestas de solución de MDVRP donde se realice la fase de asignación de clientes a depósitos usando como vía de resolución los algoritmos de agrupamiento.
  \end{itemize}
  \item Implementar en BHAVRP un nuevo algoritmo de agrupamiento basado en densidad.
  \begin{itemize}
    \item Analizar algoritmos de agrupamiento basados en densidad en cuanto a su funcionamiento para su incorporación en la nueva versión de la biblioteca.
    \item Seleccionar el algoritmo de agrupamiento que se van a incorporar a la biblioteca para resolver la fase de asignación del MDVRP.
    \item Implementar y adaptar el algoritmo seleccionado teniendo en cuenta las características que presentan.
  \end{itemize}
  \item Validar el funcionamiento del algoritmo mediante la ejecución  de experimentos.
  \begin{itemize}
    \item Realizar pruebas unitarias, donde se muestre que el algoritmo implementado funcione correctamente.
    \item Realizar pruebas estadísticas no paramétricas entre los resultados obtenidos por las heurísticas de construcción usando los resultados de los algoritmos de agrupamiento de BHAVRP.
    \item Realizar una prueba de Pearson y una prueba de Spearman para validar si existe una correlación positiva entre las métricas internas de BHAVRP y la calidad de las rutas obtenidas por las heurísticas de construcción.
  \end{itemize}
\end{enumerate}

El cumplimiento del objetivo general tiene como valor práctico un aporte significativo en la resolución de la fase de asignación del problema de planificación de rutas de vehículos con múltiples depósitos. Los resultados tienen valor científico y tecnológico ya que se pone a disposición de la comunidad científica un componente de software que brinda nuevos algoritmos para la solución de MDVRP.

Para alcanzar los objetivos mencionados, el trabajo se encuentra estructurado en tres capítulos, incluyendo secciones importantes como conclusiones, recomendaciones y referencias bibliográficas. El \autoref{cap:fundamentos}: \nameref{cap:fundamentos} donde se presenta el estado del arte del MDVRP, las estrategias de solución existentes, así como un análisis comparativo de bibliotecas. \autoref{cap:incorporacion}: \nameref{cap:incorporacion} describe los algoritmos de agrupamiento basados en densiad candidatos más comentados en la literatura, el algoritmo propuesto detallando el funcionamiento de este según la literatura y las adaptaciones al MDVRP que hay que realizar. Finalmente, el Capítulo 3: Análisis experimental expone varias pruebas unitarias de las funciones más relevantes del el algoritmo implementado, así como un estudio comparativo entre algoritmos de asignación en las heurísticas de construcción aplicados sobre ciertas instancias y un análisis de correlación entre las métricas de validación internas y la calidad de las rutas obtenidas.
%==============================================================
%                    CAPITULO 1
%==============================================================
\chapter{Fundamentos teóricos del problema de planificación de rutas y su solución mediante diferentes métodos}\label{cap:fundamentos}

En este capítulo se describen los aspectos teóricos de la investigación que se encuentran relacionados con los MDVRP. Entre los aspectos que se abordan están las características y variantes del problema, así como los diferentes métodos de solución encontrados en la literatura. Además, se hace énfasis en los algoritmos de agrupamiento como alternativa para la fase de asignación de clientes a depósitos. Por último, se describe la Biblioteca de Heurísticas de Asignación para Problemas de Planificación de Rutas de Vehículos con Múltiples Depósitos y sus algoritmos existentes.

\section{Problemas de Planificación de Rutas de Vehículos}

El concepto de Problema de Planificación de Rutas de Vehículos (\textit{Vehicle Routing Problem}, VRP) es una generalización del Problema del Viajero (\textit{Traveling Salesman Problem}, TSP). En concreto, el VRP es un problema de optimización combinatoria cuyo objetivo es atender las demandas de clientes que se encuentran dispersos geográficamente sujetos a un conjunto de restricciones. La diferencia radica en que TSP dispone de un mapa con múltiples puntos que un individuo tiene que visitar para volver a su punto de origen realizando la ruta más corta posible. En cambio, el VRP no se limita a solo una persona, sino que pueden existir más de un individuo, en este caso vehículos, que normalmente deben entregar objetos en diversos puntos de la ciudad a clientes. De este modo, el problema se vuelve más complejo que el TSP \parencite{kumarSurveyVehicleRouting2012}.

Las aplicaciones informáticas para el dise\~no de rutas se derivan de modelos VRP que atienden diferentes necesidades e incluyen diversas restricciones. Esto ha dado como resultado el surgimiento de diferentes variantes del clásico VRP, de las cuales se pueden mencionar las siguientes:

\begin{itemize}
  \item VRP con recolección y entrega (VRPPD): estudia el caso donde una empresa debe recolectar y entregar bienes en cantidades específicas para cada cliente visitado \parencite{berbegliaDynamicPickupDelivery2007}.
  \item VRP con flota heterogénea: los vehículos tienen diferentes capacidades de carga  \parencite{baldacciRoutingHeterogeneousFleet2008}.
  \item \textit{Open VRP} (OVRP): se refiere al caso de cuando el vehículo no regresa al depósito de origen, es decir, que se dirige a otro depósito o sitio diferente para cargar productos, en un contexto de flota rentada \parencite{sariklisHeuristicMethodOpen2000}.
  \item VRP con ventanas de tiempo (VRPTW): plantea que cada cliente tiene que ser atendido de manera obligada dentro de un cierto horario específico \parencite{cordeauVRPTimeWindows2002}.
  \item VRP con depósitos múltiples (MDVRP): implica que la empresa posee diversos depósitos desde los cuales puede abastecer a los clientes \parencite{montoya-torresLiteratureReviewVehicle2015}.
  \item \textit{Stochastic VRP} (SVRP): Propone que las variables involucradas son desconocidas o aleatorias, e incluso que se pueden encontrar en un rango de probabilidad, tales como el número de clientes, sus demandas, tiempo de servicio, tiempo de viaje, entre otras \parencite{gendreauStochasticVehicleRouting1996}.
  \item \textit{Periodic VRP} (PVRP): analiza un problema donde los pedidos pueden ser llevados sólo en ciertos días  \parencite{francisPeriodVehicleRouting2008}.
  \item Capacitated VRP: estudia el caso donde los vehículos tienen una capacidad máxima de carga y se deben dise\~nar rutas que satisfagan la demanda de los clientes sin exceder dicha capacidad \parencite{ibrahimCAPACITATEDVEHICLEROUTING2019}.
\end{itemize}

\subsection{Problema de Planificación de Rutas de Vehículos con Múltiples Depósitos}

El MDVRP es un problema de optimización combinatoria, categorizado por su complejidad como NP-Duro  \parencite{rebillasloredoMultidepotVRPVehicle2018}. Por lo que es difícil encontrar su solución óptima usando métodos exactos en un tiempo polinomial. Esta variante cuenta con varios depósitos y cada uno tiene su propia flota de vehículos  \parencite{colinDisenoAlgoritmoMetaheuristico2024}.

El problema puede ser representado mediante modelos de optimización y se basa en  \parencite{montoya-torresLiteratureReviewVehicle2015}, \parencite{Cordeau2012}, \parencite{Vidal2012}, \parencite{Subramanian2013}, \parencite{Escobar2014}, \parencite{cordeauTabuSearchHeuristic1997}. Sea \(G = (V, A)\) un grafo completo, donde \(V\) es el conjunto de nodos y \(A\) es el conjunto de arcos. Los nodos se dividen en dos subconjuntos: los clientes a ser atendidos, \(V_{C} = \{1, \ldots, N\}\), y los múltiples depósitos \(V_{D} = \{N + 1, \ldots, N + M\}\), con \(V_{C} \cup V_{D} = V\) y \(V_{C} \cap V_{D} = \emptyset\). Hay un costo no negativo \(c_{ij}\) asociado con cada arco \((i, j) \in A\). La demanda de cada cliente es \(d_{i}\) donde no hay demanda en los nodos de depósito. También hay una flota de \(K\) vehículos, cada uno con capacidad \(Q\). El tiempo de servicio en cada cliente \(i\) es \(t_{i}\) mientras que el tiempo máximo de duración de la ruta se establece en \(T\). Es posible que se necesite un factor de conversión \(w_{ij}\) para transformar el costo \(c_{ij}\) en unidades de tiempo. Sin embargo, el costo es el mismo que las unidades de tiempo y distancia, por lo que \(w_{ij} = 1\)  \parencite{oliveiraCooperativeCoevolutionaryModels2015}.

A continuación se presenta la formulación matemática \parencite{oliveiraCooperativeCoevolutionaryModels2015}:

Las variables binarias \(x_{ijk}\) son iguales a 1 cuando el vehículo \(k\) visita el nodo \(j\) inmediatamente después del nodo \(i\). Las variables auxiliares \(y_{i}\) también se utilizan en las restricciones de eliminación de subtours.

\noindent\textbf{Función Objetivo:} Minimizar el costo total.
\begin{equation}
\mathrm{Minimize}\sum_{i = 1}^{N + M}\sum_{j = 1}^{N + M}\sum_{k = 1}^{K}c_{ij}x_{ijk},
\end{equation}

\noindent\textbf{Restricciones:}

Garantiza que cada cliente sea visitado exactamente una vez (flujo entrante).
\begin{equation}
\sum_{i = 1}^{N + M}\sum_{k = 1}^{K}x_{ijk} = 1 \quad (j = 1,\dots ,N);
\end{equation}

Garantiza que cada cliente sea visitado exactamente una vez (flujo saliente).
\begin{equation}
\sum_{j = 1}^{N + M}\sum_{k = 1}^{K}x_{ijk} = 1 \quad (i = 1,\dots ,N);
\end{equation}

Conservación de flujo. Lo que entra a un nodo es igual a lo que sale de él para cada vehículo.
\begin{equation}
\sum_{i = 1}^{N + M}x_{ihk} - \sum_{j = 1}^{N + M}x_{hjk} = 0 \quad (k = 1,\dots ,K;h = 1,\dots ,N + M);
\end{equation}

Restricción de capacidad. La demanda total de la ruta no debe exceder la capacidad del vehículo \(Q\).
\begin{equation}
\sum_{i = 1}^{N + M}\sum_{j = 1}^{N + M}d_{i}x_{ijk}\leq Q \quad (k = 1,\dots ,K);
\end{equation}

Restricción de duración de la ruta. El costo total más el tiempo de servicio no debe exceder el tiempo máximo \(T\).
\begin{equation}
\sum_{i = 1}^{N + M}\sum_{j = 1}^{N + M}(c_{ij}w_{ij} + t_i)x_{ijk}\leq T \quad (k = 1,\dots,K);
\end{equation}

Verifica la disponibilidad de vehículos (salida del depósito).
\begin{equation}
\sum_{i = N + 1}^{N + M}\sum_{j = 1}^{N}x_{ijk}\leq 1 \quad (k = 1,\dots,K);
\end{equation}

Verifica la disponibilidad de vehículos (entrada al depósito).
\begin{equation}
\sum_{j = N + 1}^{N + M}\sum_{i = 1}^{N}x_{ijk}\leq 1 \quad (k = 1,\dots,K);
\end{equation}

Restricciones de eliminación de subtours.
\begin{equation}
y_{i} - y_{j} + (M + N)x_{ijk}\leq N + M - 1; \quad \mathrm{para} \ 1\leq i\neq j\leq N \ \mathrm{y} \ 1\leq k\leq K;
\end{equation}

Define la naturaleza binaria de las variables de decisión de ruta.
\begin{equation}
x_{ijk}\in \{0,1\} \quad \forall i,j,k;
\end{equation}

Define la naturaleza binaria de las variables auxiliares de eliminación de subtours.
\begin{equation}
y_{i}\in \{0,1\} \quad \forall i;
\end{equation}

En la Figura~\ref{fig:mdvrp} se ilustra la estructura del problema:

\begin{itemize}
  \item Se observan tres depósitos (representados por cuadrados), cada uno con una flota de vehículos.
  \item Los clientes (círculos rojos) son asignados a un depósito específico.
  \item Las líneas verdes muestran las rutas de los vehículos, que parten y regresan a su depósito de origen.
\end{itemize}

\figura{fig01_mdvrp.png}{0.75\textwidth}{Representación gráfica de MDVRP.}{fig:mdvrp}

La literatura documenta múltiples aplicaciones del MDVRP a problemas logísticos reales. Un estudio aplicado a una red de distribución de productos de limpieza y cuidado personal en Turquía utiliza el MDVRP para optimizar las rutas desde tres depósitos a 502 clientes. Mediante un enfoque de agrupar primero y planificar después se logran mejoras significativas en la eficiencia de la distribución, utilizando datos reales de demanda y coordenadas geográficas \parencite{hamidOptimizationEssentialsTheory2024}.

Otro ejemplo se encuentra en el Reino Unido, donde se desarrolla un modelo de programación entera mixta (MIP) \parencite{junger50YearsInteger2009} para resolver un MDVRP que incluye complejidades del mundo real como una flota heterogénea, costos de conductores y ventanas de tiempo. La implementación de este modelo permite reducir sustancialmente los costos de transporte sin necesidad de abrir nuevas instalaciones \parencite{mogaleModellingOptimisingMultidepot2025}.

También se encuentra un caso de estudio en Estambul que trata sobre la distribución de una empresa de bebidas utilizando datos históricos de GPS para construir una matriz de tiempo-distancia más precisa. Se evalúan múltiples enfoques de optimización, incluyendo algoritmos evolutivos y de aprendizaje profundo, para resolver el MDVRP. Los resultados muestran que estos métodos son efectivos para manejar la complejidad y la naturaleza dinámica de las entregas en entornos urbanos de alta densidad  \parencite{erdemAIInfusedEvolutionaryOptimization2025}.

\section{Algoritmos heurísticos}

Debido a la complejidad que presentan los VRP, existen dos tipos de algoritmos que pueden resolver este tipo de problemas, los métodos exactos y aproximados. Cada uno de ellos posee diferentes características que pueden considerarse ventajas y/o desventajas que deben ser tomadas en cuenta al ser utilizados \parencite{colinDisenoAlgoritmoMetaheuristico2024}. En la Figura~\ref{fig:tecnicas} se puede apreciar una taxonomía de estos métodos:

\figura{fig02_tecnicas.png}{0.85\textwidth}{Técnicas utilizadas para la solución de MDVRP \parencite{talbiMetaheuristicsDesignImplementation2009}.}{fig:tecnicas}

Entre los métodos exactos se pueden distinguir tres grupos: búsqueda directa de árbol, programación dinámica y programación lineal, entera y mixta. Todos estos métodos son eficientes en problemas de hasta 50 clientes, debido a restricciones de tiempo computacional. Se han llevado a cabo algunas investigaciones introduciendo métodos exactos para MDVRP  \parencite{colinDisenoAlgoritmoMetaheuristico2024}. Algunos ejemplos de los más utilizados para el MDVRP de este tipo de métodos son:

\begin{itemize}
  \item Ramificación y Corte: los algoritmos \textit{Branch \& Cut} se basan en el cálculo simultáneo de cotas inferiores y superiores de la solución óptima. Las cotas superiores se obtienen por medio de la resolución de las relajaciones lineales asociadas a cada nodo del árbol de enumeración, eventualmente ajustados por medio de planos de corte, mientras que las cotas inferiores son halladas por medio de procedimientos heurísticos y cuando la relajación lineal asociada con algún subproblema tiene óptimo entero  \parencite{hernandezAPLICACIONMODELOVRP2015}.
  \item Programación Lineal: la programación lineal es un método a través del cual se optimiza una función objetivo, bien sea maximizando o minimizando dicha función, en la cual las variables están elevadas a la potencia 1. Este campo de la programación lineal es un área de la programación matemática, dedicada a maximizar o minimizar una función lineal que recibe el nombre de ``función objetivo'', de manera que las variables de tal función estén sujetas a una serie de restricciones expresadas a través de un sistema de ecuaciones o inecuaciones lineales \parencite{colinDisenoAlgoritmoMetaheuristico2024}.
\end{itemize}

La investigación se centra en los métodos heurísticos, los cuales se han caracterizado por ser procedimientos que garantizan buenas soluciones mediante una exploración limitada del espacio de búsqueda. Los métodos heurísticos se pueden clasificar en métodos constructivos y heurísticas de mejora \parencite{colinDisenoAlgoritmoMetaheuristico2024}. Algunas heurísticas utilizadas en la literatura para solucionar el MDVRP son:

\begin{itemize}
  \item Heurística ad hoc: en 1993 fue presentada una nueva heurística ad hoc en el aporte de \parencite{chaoNewHeuristicMultidepot1993}, su método de solución asigna los clientes a los depósitos de manera fácil y rápida. En el procedimiento de mejora, permite que la distancia total recorrida por la flota aumente con la esperanza de encontrar finalmente rutas muy cortas \parencite{colinDisenoAlgoritmoMetaheuristico2024}.
  \item Algoritmo de Ahorros: extensión del método clásico para VRP al contexto de múltiples depósitos. Asigna clientes a depósitos y luego construye rutas fusionando recorridos según el ahorro de distancia \parencite{clarkeSchedulingVehiclesCentral1964}. En  \parencite{tillmanMultipleTerminalDelivery1969} se presentó la primera adaptación para MDVRP.
  \item Algoritmo de Barrido: método de dos fases: primero agrupa clientes por ángulo polar respecto al depósito, luego construye rutas para cada grupo. Bueno para problemas de tama\~no medio-grande \parencite{gillettHeuristicAlgorithmVehicledispatch1974}.
  \item \textit{Merge-Head and Drop-Tail}: una técnica planteada para resolver el MDVRP es la denominada estrategia \textit{Merge-Head and Drop-Tail} (MHDT), combinada con optimización de combinación de sectores (SCO), aplicada por  \parencite{shiHeuristicSolutionMethod2020}; aquí se plantea un comportamiento inverso a la entrega, es decir debe generar rutas de recolección, dada esta única diferencia, el resto de las características y restricciones son idénticas, su propuesta funciona de igual manera para el MDVRP. En este método, primero asignan puntos de recolección de acuerdo con la distancia más cercana a los puntos de recolección, luego cada planta resuelve el VRP respectivamente, asignando clientes a los vehículos y planificando el orden en que los vehículos visitan a los clientes. En el último paso, proponen el algoritmo de optimización de combinación de sectores (SCO) para generar múltiples soluciones iniciales, y luego estas soluciones iniciales se mejoran utilizando la estrategia de combinación de cabeza y caída (MHDT).
\end{itemize}

Las metaheurísticas son estrategias de alto nivel que guían y modifican heurísticas subyacentes para explorar el espacio de búsqueda de forma eficiente. Estos algoritmos se basan en la exploración iterativa del espacio de soluciones, utilizando estrategias para escapar de los óptimos locales. Estas se clasifican en basadas en población o basadas en un punto  \parencite{blumMetaheuristicsCombinatorialOptimization2003}.

La Búsqueda Tabú se mantiene como uno de los mejores enfoques para resolver el MDVRP. Esta es una metaheurística que explora el espacio de soluciones permitiendo movimientos que empeoran la solución actual para escapar de óptimos locales. Para evitar ciclos, mantiene una lista de movimientos prohibidos o tabú durante un número determinado de iteraciones  \parencite{hernandezAPLICACIONMODELOVRP2015}.

También otra de las metaheurísticas de punto único que se destaca es la Búsqueda Local Iterada. Este algoritmo parte de una solución inicial, aplica una perturbación para diversificar la búsqueda, luego una búsqueda local para alcanzar un nuevo óptimo y finalmente decide si acepta la nueva solución según algún criterio \parencite{colinDisenoAlgoritmoMetaheuristico2024}.

Entre los poblacionales se destacan bastante los genéticos por la flexibilidad ante ciertas restricciones. Estos algoritmos mantienen una población de soluciones que evoluciona mediante operadores de selección, cruce y mutación. Para el MDVRP, las soluciones se representan como secuencias de clientes. El cruce combina sub-rutas de dos padres para generar un descendiente, reparando la factibilidad si es necesario \parencite{hernandezAPLICACIONMODELOVRP2015}.

Otro de los algoritmos poblacionales que se destacan fue el Algoritmo Colonia de Hormigas. Esta simula el comportamiento de hormigas reales que depositan feromona en el camino. Las hormigas artificiales construyen soluciones probabilísticamente guiadas por la feromona (que refleja la calidad de soluciones previas) y por información heurística local. Para MDVRP, se optimizan rutas por depósito y se a\~naden operadores de cruce para mejorar el rendimiento \parencite{colinDisenoAlgoritmoMetaheuristico2024}.

\section{Algoritmos de agrupamiento para la fase de asignación de clientes a depósitos}

La minería de datos es un campo interdisciplinario de la computación que se centra en el descubrimiento de patrones significativos en grandes volúmenes de datos mediante el uso de métodos avanzados de inteligencia artificial, aprendizaje automático, estadística y bases de datos. Su objetivo principal es extraer conocimiento valioso de los datos y transformarlo en una estructura comprensible, ya sea para su análisis o para su almacenamiento en bases de datos con miras a usos posteriores \parencite{oralloIntroduccionMineriaDatos2004}.

La minería de datos se basa en una variedad de técnicas y métodos que permiten extraer información útil de grandes volúmenes de datos. Algunas de las técnicas más comunes incluyen \parencite{kesavarajStudyClassificationTechniques2013}:

\begin{itemize}
  \item Clasificación: es una técnica utilizada para asignar elementos de un conjunto de datos a categorías predefinidas. Esto se logra mediante algoritmos de aprendizaje supervisado, como los árboles de decisión, las máquinas de vectores de soporte (SVM) y las redes neuronales.
  \item Agrupamiento: es una técnica no supervisada que organiza los datos en clústeres basados en la similitud entre los datos. Esta no requiere categorías predefinidas.
  \item Regresión: es una técnica utilizada para predecir un valor numérico basado en la relación entre variables.
  \item Reglas de asociación: identifican relaciones entre variables en un conjunto de datos.
\end{itemize}

La fase de asignación de clientes a depósitos en MDVRP puede considerarse un problema de agrupamiento, ya que el resultado es un conjunto de clientes agrupados por depósito. Esta interpretación sugiere un enfoque que agrupa a los clientes y luego programa los vehículos en cada grupo. En la literatura el MDVRP puede resolverse en dos etapas: primero, los clientes deben asignarse a los depósitos; luego, se deben construir rutas que vinculen a los clientes asignados al mismo depósito. Un enfoque razonable sería dividir el problema en tantos subproblemas como depósitos existan y resolver cada subproblema por separado, lo cual transforma el MDVRP en múltiples VRP independientes \parencite{sekerSolvingMultiDepot2009}. Por tanto, resulta interesante adaptar los algoritmos de agrupamiento para la fase de asignación en el MDVRP.

Los algoritmos de agrupamiento se basan en la cuantificación de qué tan similares o no sean los objetos del conjunto de datos. Las métricas de distancia desempe\~nan un papel central en la agrupación de estos objetos en grupos. Entre las medidas comunes de similitud y distancia se incluyen varios tipos de distancia como \textit{Euclidian}, \textit{Manhattan}, \textit{Haversine} y \textit{Chebyshev}. La elección de la medida de proximidad influye en los resultados de la agrupación, ya que los distintos conjuntos de datos requieren técnicas de separación específicas y métricas de similitud adaptadas a sus características. Estos algoritmos de agrupamiento según  \parencite{koutroumbasPatternRecognition2008}, se pueden clasificar como:

\begin{itemize}
  \item Agrupamiento particional: en estos algoritmos cada partición representa un grupo y cada objeto debe pertenecer exactamente a un grupo. Si se aplica la partición difusa, permite que los objetos pertenezcan a múltiples grupos con probabilidades. Aquí se encuentran algoritmos como k-means, k-medoids, k-modes y c-means difusos. Estos métodos son sensibles a la inicialización y a los valores atípicos, y se obtienen resultados confiables cuando los grupos son esféricos, de tama\~no y distribución uniformes. En la Figura~\ref{fig:particional} se puede apreciar gráficamente el funcionamiento de este tipo de agrupamientos.
  \item Agrupamiento jerárquico: se dividen en enfoques basados en aglomeración (de abajo hacia arriba) y divisivos (de arriba hacia abajo), que producen una jerarquía de grupos representados como dendrogramas. Los métodos basados en aglomeración comienzan con cada punto de datos como un grupo separado y fusionan los grupos más cercanos de forma iterativa, mientras que los métodos divisivos comienzan con todos los puntos de datos en un grupo y los dividen de forma recursiva. Los criterios de vinculación incluyen la vinculación simple, completa, promedio y el método de Ward, cada uno de los cuales define la proximidad del grupo de manera diferente. Los ejemplos de algoritmos de agrupamiento jerárquico incluyen BIRCH, CURE y ROCK. En la Figura~\ref{fig:jerarquico} se puede apreciar gráficamente el funcionamiento de este tipo de agrupamientos.
  \item Agrupamiento basado en densidad: los algoritmos de agrupamiento basados en densidad identifican los agrupamientos como regiones de alta densidad separadas por áreas de baja densidad, sin requerir el número de agrupamientos como entrada. Dentro de esta clasificación se encuentra el agrupamiento espacial basado en densidad de aplicaciones con ruido (DBSCAN), el agrupamiento de puntos para identificar la estructura del agrupamiento (OPTICS), el agrupamiento de picos de densidad (DPC) y el agrupamiento jerárquico basado en densidad (HDBSCAN), este último destaca por construir una jerarquía de clústeres que le permite manejar densidades variables sin requerir un parámetro global de radio. Estos métodos pueden presentar dificultades con densidades variables. En la Figura~\ref{fig:densidad} se puede apreciar gráficamente el funcionamiento de este tipo de agrupamientos.
\end{itemize}

\figura{fig03_particional.png}{1\textwidth}{Representación gráfica de algoritmos de agrupamiento particionales \parencite{gaoOverviewClusteringMethods2023}.}{fig:particional}

\figura{fig04_jerarquico.png}{1\textwidth}{Representación gráfica de algoritmos de agrupamiento jerárquicos \parencite{gaoOverviewClusteringMethods2023}.}{fig:jerarquico}

\figura{fig05_densidad.png}{1\textwidth}{Representación gráfica de algoritmos de agrupamiento basados en densidad \parencite{gaoOverviewClusteringMethods2023}.}{fig:densidad}

Los algoritmos de agrupamiento se aplican al MDVRP principalmente bajo el enfoque ``agrupar primero, planificar después'', donde primero se agrupan clientes y luego se construyen rutas para cada grupo. Un buen ejemplo de la aplicación de esto es una distribución de alimentos a supermercados en Reino Unido donde se utilizó K-means para agrupar 120 supermercados en cuatro clústeres correspondientes a los depósitos disponibles, minimizando la distancia intra-cluster. Luego se optimizaron rutas independientes para cada clúster. El ahorro en costos operativos fue del 15\% \parencite{salhiHeuristicAlgorithmMultidepot1998}.

\section{Biblioteca de Heurísticas de Asignación en la Planificación de Rutas de Vehículos con Múltiples Depósitos}

El área de optimización ha llevado a la creación de diversas bibliotecas en diferentes lenguajes de programación, cada una con su propia variedad de algoritmos y características. A continuación, se presenta una lista no exhaustiva de bibliotecas destacadas en resolver VRP y sus variantes, entre las cuales se incluye el MDVRP. De cada una se ofrece una breve descripción de sus características y usos  \parencite{ericramosaragonIncorporacionModuloMetricas2025}:

\begin{itemize}
  \item JSprit: es una biblioteca de código abierto escrita en Java. Es un componente ligero y flexible que resulta fácil de utilizar, fundamentado en una única metaheurística multipropósito que actualmente aborda una variedad amplia de problemas, como el Problema del Agente Viajero y sus variantes de VRP. Aunque no está específicamente dise\~nada para abordar la fase de asignación en problemas con múltiples depósitos, JSprit puede implementar estrategias de agrupamiento geográfico como un paso preliminar, utilizando técnicas como ordenamientos angulares o agrupaciones espaciales similares al enfoque del clásico algoritmo de barrido. Estas estrategias permiten organizar a los clientes en función de su proximidad antes de aplicar su metaheurística interna  \parencite{jspritcontributorsJSpritJavabasedLibrary2026}.
  \item \textit{Google Operation Research Tools} (OR-Tools): es un software de código abierto desarrollado en C++, con versiones estables disponibles en Java, C\# y Python. Originalmente dise\~nado para abordar diversos problemas de optimización como la programación lineal, la programación entera y la programación con restricciones, esta herramienta ha ido incorporando gradualmente capacidades para resolver variantes del VRP. Aunque brinda herramientas para construir modelos completos de VRP, su arquitectura flexible permite personalizar fases específicas, como la asignación, mediante programación y heurísticas personalizadas. Además, se pueden integrar algoritmos de agrupamiento para realizar una asignación previa de clientes a los depósitos. También es posible implementar restricciones a medida en el modelo, como límites de capacidad o balanceo de carga, para ajustar la asignación inicial antes de continuar con la planificación de las rutas de los vehículos. Esta herramienta resuelve varias variantes de VRP de manera nativa, pero no el MDVRP \parencite{perronORTools}.
  \item \textit{Vehicle Routing Open-Source Optimization Machine} (VROOM): es un componente de software para problemas de planificación de rutas de vehículos desarrollado en el lenguaje de programación C++. Esta herramienta utiliza \textit{Open Street Map} (OSM) y \textit{Open Source Routing Machine} (OSRM) como \textit{backend} para obtener rutas y devolver soluciones para las diferentes variantes del VRP. VROOM puede resolver distintas variantes del VRP que incluyen múltiples depósitos, capacidades, ventanas de tiempo, así como entrega y recogida \parencite{karkulaComparisonCapabilitiesRecent2019}.
  \item Scikit-learn: es una biblioteca de aprendizaje automático en Python que, aunque no está dise\~nada específicamente para resolver problemas de optimización como el MDVRP, proporciona una amplia gama de algoritmos de agrupamiento y análisis de datos que pueden ser utilizados para abordar la fase de asignación en este tipo de problemas. Algoritmos como K-means, DBSCAN y el agrupamiento jerárquico son útiles para preprocesar los datos y asignar clientes a depósitos de manera eficiente, basándose en proximidad y restricciones de capacidad. Su flexibilidad y compatibilidad con otras bibliotecas de Python la convierten en una herramienta versátil para integrar técnicas de agrupamiento en soluciones personalizadas para el MDVRP  \parencite{f.pedregosaScikitlearnMachineLearning2011}.
\end{itemize}

En el ámbito nacional, se han desarrollado diversas soluciones en el contexto académico orientadas a la resolución de problemas de planificación de rutas de vehículos, como parte del proyecto de Optimización y Metaheurísticas de la Facultad de Ingeniería Informática de la Universidad Tecnológica de La Habana ``José Antonio Echeverría'' (CUJAE). Algunas de estas herramientas se caracterizan por su enfoque modular y por estar dise\~nadas para abordar distintas variantes del VRP, con especial énfasis en la utilización de heurísticas y metaheurísticas aplicadas al MDVRP \parencite{ericramosaragonIncorporacionModuloMetricas2025}:

\begin{itemize}
  \item Bibliotecas de Heurísticas de Construcción para Problemas de Planificación de Rutas de Vehículos (BHCVRP): es un componente que implementa diez heurísticas de construcción para resolver siete variantes del VRP, incluyendo el MDVRP. Actualmente las soluciones obtenidas mediante estas heurísticas pueden ser utilizadas como solución final o como parte de una solución más compleja en un contexto donde se empleen metaheurísticas \parencite{anandadelacaridadmoralesmoralesNuevaVersionBiblioteca2025}.
  \item Capa de servicio para soluciones a Problemas de Planificación de Rutas de Vehículos (CaSVRP): es una capa de servicios web dirigidos a la resolución de diferentes variantes de VRP mediante algoritmos metaheurísticos. Entre las variantes del VRP que resuelve se halla MDVRP \parencite{g.yanezDesarrolloNuevosModelos2021}. Esta solución se compone de un módulo de optimización donde se utiliza la Biblioteca de Clases para la Integración de Algoritmos Metaheurísticos (BiCIAM)\parencite{hernandezNuevaVersionBiblioteca2014a}, un módulo para la generación de soluciones iniciales empleando BHCVRP y otro para la generación de nuevas soluciones a partir de otras soluciones haciendo uso de LMOVRP \parencite{loresOperadoresMutacionBasados2017}.
  \item Biblioteca de Heurísticas de Asignación para Problemas de Planificación de Rutas de Vehículos con Múltiples Depósitos (BHAVRP): es una biblioteca que resuelve la fase de asignación de clientes a depósitos en el MDVRP. Esta herramienta contiene 20 algoritmos de agrupamiento y heurísticas clásicas de asignación, y es utilizada por otras herramientas como BHCVRP, CaSVRP y TransO \parencite{ericramosaragonIncorporacionModuloMetricas2025}.
\end{itemize}

La Tabla~\ref{tab:bibliotecas} presenta una comparación entre las bibliotecas antes mencionadas que resuelven la fase de asignación del MDVRP. Es importante resaltar criterios como el hecho de ser de código abierto ya que eso sería de utilidad para seguir mejorando el componente.

\begin{table}[H]
\centering
\caption{Comparación entre las herramientas que resuelven el MDVRP.}
\label{tab:bibliotecas}
\small
\begin{tabular}{@{}lcccc@{}}
\toprule
\textbf{Biblioteca} & \textbf{\shortstack{Implementan\\heurísticas clásicas}} & \textbf{\shortstack{Implementan\\algoritmos de\\agrupamiento}} & \textbf{\shortstack{Lenguajes\\soportados}} & \textbf{\shortstack{Código\\abierto}} \\
\midrule
JSprit & No & No & Java & Sí \\
OR-Tools & Sí & No & \shortstack{C++, Python,\\Java, C\#} & Sí \\
VROOM & No & No & C++ & No \\
Scikit-learn & No & Sí & Python & Sí \\
CaSVRP & Sí & No & Java & Sí \\
BHCVRP & Sí & No & \shortstack{Java, Python} & Sí \\
BHAVRP & Sí & Sí & \shortstack{Java, Python} & Sí \\
\bottomrule
\end{tabular}
\end{table}

A partir del análisis comparativo de las bibliotecas podemos ver que solo BHAVRP y Scikit-learn ofrecen algoritmos de agrupamiento, lo que los posiciona como herramientas especialmente útiles cuando las necesidades incluyen esta capacidad para el preprocesamiento o modelado del problema. BHAVRP es la única biblioteca específicamente enfocada en resolver la fase de asignación del MDVRP, mediante la implementación de heurísticas clásicas y algoritmos de agrupamiento. Las demás bibliotecas analizadas abordan el problema completo o incluyen la asignación como un paso preliminar, sin especialización explícita en esta etapa.

BHAVRP actualmente ha sido integrado como parte funcional en herramientas como: CaSVRP y BHCVRP. Aunque su uso se ha concentrado en estos sistemas, su arquitectura modular y sus implementaciones en Java y Python permiten su potencial integración en otras aplicaciones compatibles con estos entornos \parencite{ericramosaragonIncorporacionModuloMetricas2025}.

Como se puede apreciar en la Figura~\ref{fig:flujo}, el flujo de ejecución comienza en el software que consume BHAVRP, donde se introducen los datos correspondientes a una instancia del problema MDVRP. Estos datos se envían al componente BHAVRP, el cual procede a cargar los datos del problema y a determinar el tipo de distancia que se utilizará en el cálculo de costos. A continuación, BHAVRP verifica si los datos fueron cargados correctamente. Si los datos presentan errores, el proceso se detiene. En caso contrario, se configuran los parámetros de ejecución necesarios, y se genera la matriz de costo correspondiente según la distancia seleccionada. Una vez generada la matriz de costos, el algoritmo se ejecuta en el contexto de BHAVRP, utilizando los parámetros definidos y la matriz de costos generada previamente. Al finalizar la ejecución del algoritmo, BHAVRP genera una solución, que es devuelta al sistema solicitante. Este último obtiene así los resultados de la asignación de clientes a depósitos, completando el ciclo de interacción  \parencite{ericramosaragonIncorporacionModuloMetricas2025}.

\figura{fig06_flujo_bhavrp.png}{1\textwidth}{Diagrama de actividades del flujo de BHAVRP.}{fig:flujo}

Para soportar la variedad de algoritmos y garantizar su correcta ejecución, BHAVRP presenta una arquitectura n-capas con un enfoque basado en reutilización, promoviendo consistencia con los lenguajes de desarrollo seleccionados  \parencite{ericramosaragonIncorporacionModuloMetricas2025}. Esta arquitectura se presenta en la Figura~\ref{fig:arquitectura}, que ofrece una vista detallada de BHAVRP.

\figura{fig07_arquitectura.png}{1\textwidth}{Diagrama de arquitectura n-capas con enfoque en reutilización de BHAVRP.}{fig:arquitectura}

La capa específica contiene los elementos responsables de gestionar la ejecución del sistema, como la clase Controller, que coordina la carga de datos, la configuración de parámetros y el flujo general de ejecución. Por su parte, en la capa general, que representa el núcleo de la biblioteca es responsable de la definición del problema, la modelación de las instancias y la ejecución de los algoritmos de asignación. Por su parte, la biblioteca trabaja con una estructura de datos que representan los elementos clave del MDVRP: clientes, depósitos, y sus respectivas restricciones y parámetros. La estructura básica para utilizar la biblioteca incluye un encabezado con el número de depósitos, clientes y vehículos disponibles por depósito, seguido de las capacidades máximas de los vehículos. Los clientes y los depósitos están definidos por un identificador y sus coordenadas (x, y). Además, cada cliente tiene una demanda asociada, mientras que cada depósito cuenta con una capacidad equivalente a la suma de las capacidades de los vehículos en su flota. BHAVRP en la capa intermedia contiene varias herramientas. Una de estas permite trabajar tanto con distancias reales como aproximadas (Euclidean, Manhattan, Chebyshev y Haversine). Luego para abordar el cálculo aproximado de las distancias, integra la herramienta OSRM (Open Source Routing Machine), que permite calcular distancias y tiempos de viaje basados en mapas reales, tanto en modalidad local como a través de solicitudes a un servidor en línea. En la capa de software se implementa la comunicación de la herramienta OSRM que expone sus servicios a través de una API HTTP RESTful. El cliente Java (como osrm-java-client) construye las peticiones HTTP y convierte las respuestas a JSON \parencite{ericramosaragonIncorporacionModuloMetricas2025}.

\subsection{El problema del ruido en los algoritmos de BHAVRP}
\label{sec:ruido_mdvrp}

La fase de asignación de clientes a depósitos en el MDVRP se ha abordado dentro de la biblioteca principalmente mediante heurísticas clásicas de asignación y algoritmos de agrupamiento particionales. Sin embargo, estos métodos presentan la limitación de que no están diseñados para detectar ni aislar puntos atípicos. Como resultado, generan clústeres que incluyen clientes alejados del depósito asignado lo que incrementa el costo de las rutas posteriores.

Esta limitación se evidencia en los resultados experimentales reportados en \parencite{ericramosaragonIncorporacionModuloMetricas2025} sobre una instancia real de 154 clientes y 5 depósitos en La Habana. La Figura~\ref{fig:ruido_tcc} muestra la asignación generada por la heurística \textit{Three Criteria Clustering} (TCC), donde se observa que algunos clientes asignados a un depósito se encuentran en zonas alejadas de su depósito principal. De forma similar, la Figura~\ref{fig:ruido_ff} presenta la asignación del algoritmo \textit{Farthest-First} (FF), caracterizada por clústeres con más dispersión y solapamiento entre zonas de distintos depósitos.

\figura{fig27_tcc.png}{0.85\textwidth}{Asignación generada por el algoritmo TCC sobre una instancia real.}{fig:ruido_tcc}

\figura{fig28_ff.png}{0.85\textwidth}{Asignación generada por el algoritmo FF sobre una instancia real.}{fig:ruido_ff}

Esto se ve reflejado en las métricas internas de validación. Algoritmos con más ruido obtienen los peores valores en métricas de cohesión y separación, lo que confirma que la presencia de este degrada la calidad estructural de la partición \parencite{ericramosaragonIncorporacionModuloMetricas2025}. Adicionalmente, el análisis de correlación entre las métricas internas y el costo total de las rutas evidencia que una asignación deficiente no siempre se traduce en un costo mayor.

Estos resultados fundamentan la necesidad de incorporar a BHAVRP un algoritmo de agrupamiento basado en densidad, capaz de identificar regiones de alta concentración de clientes y aislar aquellos puntos que no pertenecen a ninguna estructura densa. A diferencia de los métodos particionales y las heurísticas clásicas, los algoritmos basados en densidad como DBSCAN, HDBSCAN u OPTICS incorporan explícitamente un mecanismo de detección de ruido, lo que permite tratarlos de forma diferenciada en la fase de asignación. Aunque la biblioteca ya contiene una adaptación de DBSCAN este presenta ciertas limitaciones a considerar, por lo que es necesario implementar otro algoritmo de esta categoría.

\section{Conclusiones parciales}

A partir del estudio realizado sobre el marco teórico de los problemas de planificación de rutas de vehículos, específicamente la variante con múltiples depósitos, así como los métodos de solución y algoritmos de Minería de Datos existentes para resolver estos problemas se puede concluir que:

\begin{itemize}
  \item Los VRP son problemas de optimización combinatoria clasificados como NP-Duro por lo que se utilizan principalmente algoritmos heurísticos y metaheurísticos para su solución. El MDVRP es una de las numerosas variantes del VRP y requiere otro tipo de estrategias de solución.
  \item Una de las estrategias para la solución del MDVRP es dividir en la fase de asignación y la fase de planificación de las rutas. Además, para esta primera fase de asignación resulta efectivo el uso de algoritmos de minería de datos.
  \item BHAVRP es una biblioteca que permite resolver esa primera fase de asignación y tiene la facilidad para a\~nadir nuevas heurísticas y algoritmos de agrupamiento gracias a su arquitectura n-capas orientada a reutilización.
  \item La incorporación de otro algoritmo basado en densidad a BHAVRP enriquece la biblioteca con otro algoritmo de esta categoría que permita tener un mecanismo de detección de ruido.
\end{itemize}
%==============================================================
%                    CAPITULO 2
%==============================================================
\chapter{Incorporaci\'on de algoritmo a la Biblioteca de Heur\'isticas de Asignaci\'on del Problema de Planificaci\'on de Rutas con M\'ultiples Dep\'ositos}\label{cap:incorporacion}

En este cap\'itulo se presenta la incorporaci\'on del nuevos algoritmo en BHAVRP. Se muestra una tabla comparativa entre los algoritmos de agrupamiento candidatos para implementar en la biblioteca. Adem\'as, se presenta una descripci\'on detallada de el algoritmo de agrupamiento elegido, modificaciones en la arquitectura de la biblioteca para la incorporaci\'on de este y la adaptaci\'on para el MDVRP mediante diagrama de actividades.

\section{Análisis y selección del algoritmo basado en densidad}

La incorporación de un algoritmo basado en densidad responde a una necesidad identificada en  la biblioteca BHAVRP carece de una variedad suficiente de algoritmos de esta categoría para garantizar un mecanismo deteccción de ruido en las asignaciones de clientes de depósitos. Las heuristicas clásicas y los algoritmos particionales son sensibles al ruido, mientras que los jerárquicos son más flexibles pero pueden verse afectados por el ruido \parencite{jainAlgorithmsClusteringData1988}. Sin embargo, los algoritmos de agrupamiento basados en densidad identifican clústeres como regiones de alta densidad separadas por áreas de baja densidad siendo capaces de detectar formas arbitrarias e identificar correctamente el ruido \parencite{esterDensitybasedAlgorithmDiscovering1996}.

Como candidatos a incluir en la biblioteca se seleccionaron los algoritmos de agrupamiento basados en densidad más comentados en la literatura, los cuales se presentan a continuación:

\begin{itemize}
\item OPTICS: genera un ordenamiento de los puntos basado en su densidad de alcanzabilidad, lo que permite identificar estructuras jerárquicas y densidades variables sin requerir un parámetro global de radio. Sin embargo, no devuelve directamente una partición de clústeres, sino que requiere un paso posterior de extracción. Es uno de los algoritmos de referencia más citados en la literatura de agrupamiento basado en densidad \parencite{ankerst1999optics}.

\item HDBSCAN: es una versión extendida de DBSCAN que construye un árbol jerárquico de clústeres basado en la distancia de alcanzabilidad mutua. Extrae los clústeres más estables, lo que le permite manejar densidades variables, detectar formas arbitrarias y aislar el ruido de forma robusta, manteniendo una complejidad computacional adecuada para grandes volúmenes de datos. Está considerado como el estado del arte dentro de los algoritmos jerárquicos basados en densidad por su enfoque híbrido de ambos enfoques \parencite{campello2013density}.

\item DENCLUE: modela el espacio de datos mediante funciones de influencia (como gaussianas) para estimar la densidad. Los clústeres se identifican como los máximos locales de esta función. Su rendimiento depende en gran medida de la correcta configuración del ancho de banda y del umbral de densidad. Se considera uno de los métodos pioneros en combinar funciones de densidad con técnicas de agrupamiento \parencite{hinneburg1998efficient}.

\item DPC: identifica los centros de los clústeres buscando puntos con alta densidad local y una distancia considerable a puntos de mayor densidad, para luego asignar el resto de los puntos. Aunque es conceptualmente simple y no requiere especificar el número de clústeres, su complejidad computacional es \(O(n^2)\) y carece de un mecanismo explícito para el etiquetado de ruido. Es uno de los algoritmos de agrupamiento por densidad más populares de la última década, con miles de citas y múltiples extensiones \parencite{rodriguez2014clustering}.
\end{itemize}

Para fundamentar la selección del algoritmo que se incluirá en BHAVRP, en la Tabla~\ref{tab:comparate} se realizó un análisis comparativo por sus caracteristicas mas relevantes para resolver la fase de asignación del MDVRP:

\clearpage
\begin{table}[htbp]
\centering
\caption{Tabla comparativa de algoritmos de agrupamiento basados en densidad.}
\label{tab:comparate}
 \begin{tabular}{@{}lcccc@{}}
\toprule
\textbf{Caracterstica} & \textbf{{OPTICS}} & \textbf{{HDBSCAN}} & \textbf{{DENCLUE}} & \textbf{{DPC}} \\
\midrule
Devuelve cluster como resultado & No    & Sí    & Sí    & Sí \\
Detección de Ruido & Sí    & Sí    & Sí    & Parcial \\
Pocos Parámetros & Sí    & Sí    & No    & No \\
Escalabillidad & Media    & Alta    & Media    & Baja \\
Maneja densidades variables & Parcial    & Sí    & Sí    & No \\
Forma de los clústeres detectados & Arbitraria & Arbitraria & Arbitraria & Limitada \\
\bottomrule
\end{tabular}%
\end{table}%

A partir del análisis de la Tabla~\ref{tab:comparate}, se seleccionó HDBSCAN como el algoritmo a implementar. Los criterios que fundamentan esta decisión son los siguientes:

\begin{itemize}
    \item HDBSCAN es el único algoritmo que cumple con todos los criterios deseables de forma positiva. Devuelve clústeres como resultado, detecta ruido, requiere pocos parámetros, posee una escalabilidad alta, maneja densidades variables y detecta formas arbitrarias en los clústeres.
    
    \item Aunque OPTICS detecta ruido y requiere pocos parámetros, no devuelve clústeres como resultado directo que es un requisito muy importante para la fase de asignación del MDVRP, solo maneja densidades variables de forma parcial y su escalabilidad es media.
    
    \item Si bien DENCLUE devuelve clústeres, detecta ruido y maneja densidades variables, presenta dos desventajas clave frente a HDBSCAN, requiere muchos parámetros lo cual dificulta la selección de un parámetro global dentro de la biblioteca que permita un buen funcionamiento para todas las instancias posibles y su escalabilidad es solo media.
    
    \item El algoritmo DPC es el menos adecuado según los criterios evaluados. Requiere demasiados parámetros, su escalabilidad es baja, no maneja densidades variables, la forma de los clústeres que detecta es limitada y su detección de ruido es solo parcial.
\end{itemize}

Dado que BHAVRP ya contaba con una implementación de DBSCAN en su versión anterior, resulta interesante comparar la implementación existente con la nueva propuesta para entender las ventajas que HDBSCAN aporta a la biblioteca. La Tabla~\ref{tab:dbscan_vs_hdbscan} resume las diferencias clave.

\begin{table}[H]
\centering
\caption{Comparación entre DBSCAN y HDBSCAN.}
\label{tab:dbscan_vs_hdbscan}
\small
\begin{tabular}{@{}lcc@{}}
\toprule
\textbf{Característica} & \textbf{DBSCAN} & \textbf{HDBSCAN} \\
\midrule
Parámetro \(\varepsilon\) (radio) & Requerido & No requerido \\
Manejo de densidades variables & No & Sí \\
Extracción de clústeres & Por umbral global & Por estabilidad jerárquica \\
Detección de ruido & Sí & Sí \\
Complejidad computacional & \(O(n \log n)\) & \(O(n^2)\) \\
Calidad de agrupamiento & Limitada por \(\varepsilon\) & Adaptativa \\
\bottomrule
\end{tabular}
\end{table}

La principal ventaja de HDBSCAN sobre DBSCAN es la eliminación del parámetro \(\varepsilon\), que en DBSCAN es extremadamente sensible y debe ajustarse manualmente para cada instancia. En su adaptación dentro de la biblioteca este se ajusta automaticamente, sin embargo esto puede fallar en ciertas instancias y empeorar la calidad de la solución. En HDBSCAN la densidad local se captura de forma adaptativa mediante la distancia de alcanzabilidad mutua, lo que permite descubrir clústeres de densidades variables sin intervención del usuario.

Sin embargo, esta mejora tiene un costo: HDBSCAN es computacionalmente más costoso que DBSCAN. Mientras que DBSCAN puede ejecutarse en \(O(n \log n)\) con un índice espacial adecuado, HDBSCAN requiere \(O(n^2)\) en su implementación actual.

La incorporación de HDBSCAN junto a DBSCAN en BHAVRP proporciona a la biblioteca dos enfoques complementarios dentro de la familia de algoritmos basados en densidad: uno muy rápido y simple con DBSCAN y otro más robusto y adaptativo con HDBSCAN. Esto enriquece las opciones disponibles para el usuario y permite tener más mecanismo para detección de ruido dentro de la biblioteca.

\section{Algoritmo HDBSCAN}

El algoritmo \textit{Hierarchical Density-Based Spatial Clustering of Applications with Noise} (HDBSCAN) fue propuesto por \parencite{campello2013density}. El algoritmo es una extensión de DBSCAN convirtiéndolo en un híbrido entre un método jerárquico y basado en densidad, que construye un dendrograma de clústeres y luego extrae una partición plana basada en la estabilidad de los mismos. Su funcionamiento se puede descomponer en cinco fases fundamentales.

\begin{enumerate}
\item Transformación del espacio: el primer paso consiste en transformar el espacio de datos para reflejar la densidad local. Para cada punto \(p\), se define su distancia al núcleo (\textit{core distance}, \(d_{core}(p)\)) como la distancia al \(k\)-ésimo vecino más cercano, donde \(k = \text{minPts}\). A partir de esta, se define la distancia de alcanzabilidad mutua (\textit{Mutual Reachability Distance}, MRD) entre dos puntos \(p\) y \(q\) como:

\begin{equation}
d_{mrd}(p, q) = \max\{ d_{core}(p), d_{core}(q), d(p, q) \}
\end{equation}

donde \(d(p, q)\) es la distancia euclidiana entre \(p\) y \(q\). Esta transformación hace que los puntos en regiones de baja densidad sean ``alejados'' entre sí, mientras que los puntos en regiones densas permanecen cercanos, lo que permite que el algoritmo sea robusto ante el ruido.

\item Construcción del árbol de expansión mínima: con las distancias MRD, se construye un grafo completo donde los nodos son los puntos de datos y el peso de cada arista \((p, q)\) es \(d_{mrd}(p, q)\). Sobre este grafo se calcula el Árbol de Expansión Mínima (\textit{Minimum Spanning Tree}, MST) \parencite{jainAlgorithmsClusteringData1988}, que conecta todos los puntos con el costo total mínimo. El MST se puede construir eficientemente utilizando el algoritmo de Prim\parencite{prim1957shortest}, con una complejidad de \(O(n^2)\) en el peor caso para espacios de baja dimensionalidad.

\item Construcción de la jerarquía de clústeres: a partir del MST, se construye una jerarquía de componentes conectados. Si se ordenan las aristas del MST por peso de menor a mayor y se eliminan progresivamente, se obtiene una secuencia de componentes que se fusionan. Este proceso genera un dendrograma que representa la estructura jerárquica de los clústeres a diferentes niveles de densidad.

\item Condensación de la jerarquía: el dendrograma completo puede contener muchas ramas que representan clústeres con muy pocos puntos. Para simplificarlo, se aplica un proceso de condensación basado en el parámetro \textit{min\_cluster\_size}. Las ramas que no alcanzan este tamaño mínimo se consideran ruido y se podan del árbol, dando como resultado un árbol condensado que solo contiene los clústeres significativos.

\item Extracción de clústeres estables: finalmente, se calcula la estabilidad de cada clúster en el árbol condensado. La estabilidad de un clúster \(C\) se define como la suma, sobre cada punto \(p \in C\), de la diferencia entre el nivel de densidad en el que el punto se une al clúster y el nivel en el que el clúster se divide. El algoritmo selecciona el conjunto de clústeres que maximiza la estabilidad total, lo que resulta en una partición robusta que no depende de un umbral de densidad fijo.
\end{enumerate}

A modo de resumen en la Figura~\ref{fig:hdbscan_act} se puede apreciar el diagrama de flujo del algoritmo HDBSCAN base.

\figura{activHDBSCAN.png}{1\textwidth}{Diagrama de flujo del algoritmo HDBSCAN.}{fig:hdbscan_act}

\section{Adaptación de HDBSCAN al contexto del MDVRP}

La aplicación directa de HDBSCAN al MDVRP no es posible, ya que el algoritmo original no contempla las restricciones de capacidad de los depósitos ni la necesidad de asignar cada clúster a un depósito específico. Por ello, se diseñó una adaptación que preserva la lógica de agrupamiento por densidad pura y añade una fase posterior de asignación heurística con restricciones de capacidad.

\subsection{Mecanismo de inicialización y configuración}

A diferencia de una ejecución estándar de HDBSCAN, donde los clústeres se forman libremente a partir de la densidad de los datos, en la adaptación propuesta el algoritmo opera exclusivamente sobre el conjunto de clientes. Los depósitos no se introducen en esta etapa, sino que su rol se reserva para la fase de reasignación posterior. 

La clase HDBSCAN hereda de AbstractDensity, la cual proporciona métodos auxiliares para el cálculo de distancias euclidianas y la gestión de clústeres. En la fase de inicialización, se cargan las listas de clientes, depósitos y se crean las estructuras de clústeres vacías a partir de los depósitos. Además, se configuran dinámicamente los dos parámetros principales del algoritmo:
\begin{itemize}
    \item \textit{minPoints} (\(m\_pts\)): Se calcula adaptativamente como \(\max(3, \ln(1.5 \cdot N))\), donde \(N\) es el número de clientes. Esto ajusta la sensibilidad a la densidad según el tamaño de la instancia.
    \item \textit{minClusterSize} (\(m\_clSize\)): Se establece en 2, garantizando que cualquier agrupación significativa tenga al menos dos clientes.
\end{itemize}

\clearpage
\subsection{Construcción de la jerarquía y extracción de clústeres}

La implementación sigue fielmente los pasos teóricos de HDBSCAN:
\begin{enumerate}
    \item Cálculo de distancias núcleo: para cada cliente, se calcula su distancia al \(m\_pts\)-ésimo vecino más cercano (\(d_{core}\)).
    \item Construcción del MST: se utiliza el algoritmo de Prim para construir el Árbol de Expansión Mínima (MST) sobre el grafo completo, donde los pesos de las aristas son las distancias de alcanzabilidad mutua (MRD), calculadas como \(\max(d_{core}(p), d_{core}(q), d(p,q))\). Para optimizar la extracción del mínimo, se emplea una estructura de datos auxiliar llamada \textit{Fibonacci Heap}, lo que reduce la complejidad de la cola de prioridad.
    \item Construcción del dendrograma: se añaden aristas de auto-conexión (self-edges) con el peso de la distancia núcleo y se procesan las aristas del MST ordenadas por peso ascendente. se utiliza la estructura llamada \textit{UnionFind} para gestionar la fusión de componentes conectados de forma eficiente, registrando el nacimiento, la muerte y la estabilidad de cada clúster candidato.
    \item Extracción por estabilidad: se calcula la estabilidad \(S(C_i)\) para cada clúster y se propaga hacia arriba en el árbol (\(\hat{S}\)). Se seleccionan aquellos clústeres que maximizan la estabilidad total, etiquetando el resto de los puntos (ruido) con el valor \(-1\).
\end{enumerate}

\subsection{Asignación de clústeres a depósitos}

Una vez que HDBSCAN ha generado los clústeres finales, se procede a la fase de asignación a los depósitos con su respectiva restricción de capacidad. El procedimiento implementado es el siguiente:

\begin{enumerate}
    \item Cálculo del centroide: para cada clúster \(C_i\), se calcula su centroide \((cx, cy)\) como el promedio de las coordenadas de sus clientes.
    \item Búsqueda del depósito más cercano: se calcula la distancia euclidiana entre el centroide del clúster y cada depósito, seleccionando el más cercano (\(d_j\)).
    \item Verificación de capacidad: se comprueba si la demanda acumulada del clúster más la demanda actual del depósito \(d_j\) no excede su capacidad total. Si se cumple, el clúster se fusiona con el depósito y se actualiza su demanda.
    \item Manejo de clústeres no asignables: si el depósito más cercano no tiene capacidad suficiente, los clientes de ese clúster se eliminan temporalmente del conjunto de clientes asignados (\textit{assignedCustomers}). Esto permite que sean procesados individualmente en la fase de finalización, evitando que un clúster completo quede sin asignar por una restricción de capacidad.
\end{enumerate}

\subsection{Manejo del ruido y reasignación individual}

Los puntos clasificados como ruido por HDBSCAN, así como los clientes de clústeres que no pudieron ser asignados por falta de capacidad, se procesan en la fase final del algoritmo. Para cada cliente no asignado se aplica la siguiente lógica:

\begin{enumerate}
    \item Se construye una cola de prioridad con los índices de los depósitos ordenados por distancia ascendente al cliente.
    \item Se recorre la cola y se asigna el cliente al primer depósito cuya capacidad restante sea suficiente para absorber su demanda individual.
    \item Si ningún depósito puede absorberlo, el cliente se añade a la lista de elementos no asignados definitivamente (\textit{unnasigned\_items}) en la solución final.
\end{enumerate}

A modo de resumen en la Figura~\ref{fig:hdbscan_act_adap} se puede apreciar el diagrama de flujo del procedimiento que realiza el algoritmo HDBSCAN adaptado para asignar los clústeres resultantes a los depósitos.

\figura{adapActivHDBSCAN.png}{1\textwidth}{Diagrama de flujo del paso final del algoritmo HDBSCAN adaptado.}{fig:hdbscan_act_adap}

Esta versión de HDBSCAN es una adaptación al MDVRP del algoritmo disponible en la librería de código abierto \textit{Java Statistical Analysis Tool} (JSAT)\parencite{raffJSATJavaStatistical}. De esta librería se reutilizaron las estructuras de datos fundamentales para garantizar el correcto funcionamiento del algoritmo. 

La primera estructura de datos importante es Fibonacci Heap, una estructura de cola de prioridad que ofrece un costo de \(O(1)\) para las operaciones de inserción y disminución de clave (\textit{decrease-key}), y \(O(\log n)\) para la extracción del mínimo \parencite{fredmanFibonacciHeapsUses1987}. Esta característica es crucial para la construcción del Árbol de Expansión Mínima (MST) mediante el algoritmo de Prim, ya que reduce su complejidad de \(O(V^2)\) a \(O(E + V \log V)\) para grafos dispersos \parencite{cormenIntroductionAlgorithms2009}. 

La segunda es la estructura \textit{Disjoint Set Union} (Union-Find), utilizada para gestionar la fusión de componentes conectados durante la construcción de la jerarquía de clústeres \parencite{gallerImprovedEquivalenceAlgorithm1964}. Gracias a las optimizaciones de compresión de caminos y unión por rango, esta estructura permite que las operaciones de búsqueda y unión tengan un costo casi constante de \(O(\alpha(n))\), donde \(\alpha\) es la función inversa de Ackermann \parencite{tarjanWorstcaseAnalysisSet1984}. 

\section{Incorporación del algoritmo propuesto a BHAVRP}

La incorporaci\'on del nuevo algoritmo de agrupamiento a la biblioteca BHAVRP requiri\'o la definici\'on de las siguientes clases:  HDBSCAN que extiende de AbstractDensity, así como un subpaquete de utilidades (\textit{utils}) que contiene las estructuras de datos auxiliares: UnionFind, FibHeap,  FibNode, Pair y Tuple3. Estas clases se integran a la arquitectura existente respetando el patr\'on de dise\~no que organiza los algoritmos de agrupamiento dentro de la biblioteca, basado en el patr\'on \textit{Template Method}. En la Tabla~\ref{tab:clases} se puede apreciar las clases con una peque\~na descripci\'on de estas:

\begin{table}[H]
\centering
\caption{Clases implementadas para la incorporaci\'on de los nuevos algoritmos de agrupamiento.}
\label{tab:clases}
\small
\begin{tabular}{@{}lp{10cm}@{}}
\toprule
\textbf{Clase} & \textbf{Descripci\'on} \\
\midrule
\textbf{HDBSCAN} & Implementa el algoritmo HDBSCAN adaptado al MDVRP. Realiza el cálculo de distancias núcleo, la construcción del MST mediante el algoritmo de Prim, la generación del dendrograma, la extracción de clústeres por estabilidad y la asignación de estos a los depósitos respetando las capacidades. \\
\textbf{UnionFind} & Implementación de la estructura de datos \textit{Disjoint Set Union} (Union-Find). Se utiliza durante el procesamiento de las aristas del MST para gestionar la fusión eficiente de componentes conectados y construir la jerarquía de clústeres (dendrograma). \\
\textbf{FibHeap} & Implementación de una cola de prioridad basada en un \textit{Fibonacci Heap}. Se utiliza en el algoritmo de Prim para extraer de forma eficiente el nodo de menor distancia (peso) durante la construcción del MST. \\
\textbf{FibNode} & Representa un nodo individual dentro del FibHeap. Almacena el valor (índice del cliente), la clave (distancia), punteros a su padre, hijo, nodos hermanos y su grado, facilitando las operaciones de consolidación y disminución de clave. \\
\textbf{Pair} & Almacena un par de valores genéricos \((X, Y)\). Se utiliza en la fase de extracción de clústeres para registrar las relaciones padre-hijo (clústeres izquierdo y derecho) en el árbol condensado. \\
\textbf{Tuple3} & Almacena una tupla de tres valores genéricos \((X, Y, Z)\). Se emplea para representar las aristas del MST, donde \(X\) e \(Y\) son los vértices conectados y \(Z\) es el peso (distancia MRD). También se usa para ordenar las aristas durante la construcción del dendrograma. \\
\bottomrule
\end{tabular}
\end{table}

Una vez definidas las clases que componen la solución, resulta pertinente analizar cómo interactúan durante la ejecución del algoritmo. La Figura~\ref{fig:aaa_act} presenta el diagrama de secuencia de la implementación base de HDBSCAN, correspondiente a la adaptación del algoritmo original de la librería JSAT. 

\figura{aaa.png}{1\textwidth}{Diagrama de secuencia del algoritmo HDBSCAN.}{fig:aaa_act}

%\clearpage
%	\begin{landscape}
%		\begin{figure}[p]
%			\centering
%			\includegraphics[width=1\paperwidth, height=1\paperwidth,keepaspectratio]{secuencia_diagrama.pdf}
%			\caption{Diagrama de secuencia de la implementación de HDBSCAN.}
%			\label{fig:secuencia_hdbscan_base}
%		\end{figure}
%	\end{landscape}
%\clearpage

Este diagrama de secuencia ilustra la interacción dinámica entre las clases descritas en la Tabla~\ref{tab:clases}, abarcando desde la inicialización hasta la extracción de clústeres por estabilidad. Es importante precisar que esta figura representa principalmente el flujo de agrupamiento por densidad.

Para esta actualizaci\'on de la biblioteca fue necesario modificar el paquete \textit{Clustering} de la capa general, ya que todas estas clases fueron implementadas en este. Como se puede apreciar en la Figura~\ref{fig:arquitectura_actualizada}, en azul se muestra el paquete modificado:

\figura{fig14_arquitectura_actualizada.png}{0.85\textwidth}{Diagrama de arquitectura actualizado.}{fig:arquitectura_actualizada}

En BHAVRP, todos los algoritmos de asignaci\'on de clientes a dep\'ositos siguen una estructura com\'un que consta de tres fases: inicializaci\'on, asignaci\'on iterativa y finalizaci\'on. Independientemente de la heur\'istica espec\'ifica, el flujo general es siempre el mismo: se preparan las estructuras de datos, se asignan los clientes a los dep\'ositos respetando las restricciones del problema, y se construye la soluci\'on final. Sin embargo, los detalles de c\'omo se realiza la inicializaci\'on, c\'omo se asignan los clientes en cada iteraci\'on y c\'omo se determina la convergencia var\'ian de un algoritmo a otro.

En este contexto, el patr\'on \textit{Template Method} se utiliza para definir el esqueleto del algoritmo en un m\'etodo plantilla que establece la secuencia invariante de ejecuci\'on. En el dise\~no de BHAVRP, la clase abstracta AssignmentTemplate (que hereda de la clase base Heuristic) encapsula el flujo del algoritmo en su m\'etodo plantilla \textit{toClustering} declarado en la interfaz IAssignment. Este m\'etodo define la secuencia de pasos invariantes:

\begin{enumerate}
  \item Inicialización (\textit{initialize}): crea las estructuras necesarias, carga los datos y configura los parámetros dinámicos de HDBSCAN.
  \item Asignación (\textit{assign}): ejecuta el algoritmo HDBSCAN completo como cálculo de distancias, MST, dendrograma, estabilidad y realiza la primera fase de asignación de clústeres a depósitos.
  \item Finalización (\textit{finish}): genera la solución final, procesando los clientes no asignados como ruido y clústeres rechazados mediante la reasignación individual.
\end{enumerate}

HDBSCAN respeta este patr\'on al implementar las operaciones abstractas definidas en la plantilla sin modificar la estructura general del flujo. El m\'etodo \textit{to\_clustering} mantiene la secuencia establecida por la clase base, mientras que cada algoritmo redefine los m\'etodos \textit{initialize}, \textit{assign} y \textit{finish} para incorporar su l\'ogica particular de agrupamiento. De esta manera, HDBSCAN implementa la expansi\'on de cl\'usteres basada en densidad y sigue el mismo esquema de ejecuci\'on definido por la biblioteca.

En la Figura~\ref{fig:diagrama_clases} se puede apreciar la jerarqu\'ia de clases donde se aplica el patr\'on Template Method y la implementaci\'on de las nuevas clases respetando este. La clase AssignmentTemplate act\'ua como clase base con el m\'etodo plantilla, y las distintas heur\'isticas de asignaci\'on sobrescriben los m\'etodos abstractos seg\'un su l\'ogica espec\'ifica. La clase HDBSCAN extiende AbstractDensity y utiliza composición con las clases auxiliares para la construcción del MST y el dendrograma.

\clearpage
	\begin{landscape}
		\begin{figure}[p]
			\centering
			\includegraphics[width=1.1\paperwidth, height=1.1\paperwidth,keepaspectratio]{fig_class_diagram.pdf}
			\caption{Diagrama de clases del paquete \textit{Clustering}.}
			\label{fig:diagrama_clases}
		\end{figure}
	\end{landscape}
\clearpage

La integración de HDBSCAN respeta los siguientes principios de diseño de software:

\begin{itemize}
  \item Principio Abierto/Cerrado: la arquitectura base de BHAVRP no requiere modificaciones para incorporar nuevos algoritmos de agrupamiento. La extensi\'on del sistema se logra mediante la creaci\'on de nuevas clases que implementan los m\'etodos definidos en la jerarqu\'ia existente.
  \item Principio de Responsabilidad Única: la clase HDBSCAN se encarga exclusivamente de la lógica de agrupamiento y asignación. Las operaciones comunes (cálculo de distancias, manejo de excepciones) se delegan a AbstractDensity, y las estructuras de datos complejas (Fibonacci Heap, Union-Find) se encapsulan en clases auxiliares independientes.
  \item Alta cohesi\'on: las operaciones relacionadas con el c\'alculo de distancias, vecindades y otras funcionalidades compartidas se agrupan en clases espec\'ificas como AbstractDensity. Esto permite que cada componente del sistema concentre funcionalidades relacionadas entre s\'i.
  \item Bajo acoplamiento: las implementaciones concretas dependen \'unicamente de las interfaces y estructuras generales proporcionadas por BHAVRP, lo que facilita la sustituci\'on o incorporaci\'on de nuevas heur\'isticas sin afectar el resto del sistema.
\end{itemize}

La complejidad computacional de HDBSCAN en su implementación estándar es \(O(n^2)\) en el peor caso \parencite{abduaziz2025minimal}, debido al cálculo de la matriz de distancias mutuas y la construcción del MST sobre un grafo completo. Sin embargo, el uso de FibHeap reduce el costo de las operaciones de la cola de prioridad. La fase de asignación de clústeres a depósitos tiene una complejidad de \(O(k \cdot m)\), donde \(k\) es el número de clústeres y \(m\) es el número de depósitos. Dado que \(k \ll n\) y \(m \ll n\), esta fase es computacionalmente despreciable en comparación con el agrupamiento.

La implementación de HDBSCAN en BHAVRP utiliza las excepciones incluidas en la biblioteca para su manejo explícito y poder para garantizar la estabilidad del algoritmo ante datos de entrada inválidos o situaciones imprevistas. Se utilizan las siguientes excepciones:

\begin{itemize}
    \item ProblemException: se lanza cuando los datos del problema son inconsistentes, por ejemplo, lista de clientes vacía, menos de dos clientes, o ausencia de depósitos.
    \item ClusterException: se lanza cuando ocurre un error durante la creación o manipulación de los clústeres.
    \item AssignmentException: encapsula errores inesperados durante la fase de asignación, incluyendo fallos en la construcción del MST o del dendrograma.
\end{itemize}

En la fase de inicialización, se realizan validaciones como:
\begin{itemize}
    \item Verificar que la lista de clientes no sea nula ni vacía.
    \item Verificar que existan al menos dos clientes para poder agrupar.
    \item Verificar que la lista de depósitos no esté vacía.
    \item Verificar que existan identificadores de depósitos para inicializar los clústeres.
\end{itemize}

Si alguna validación falla, se lanza la excepción correspondiente con un mensaje descriptivo, permitiendo que el sistema que consume BHAVRP pueda reaccionar adecuadamente y buscar una solución con mayor facilidad.

\clearpage

\section{Conclusiones parciales}

A partir del análisis, diseño e implementación descritos en este capítulo, se pueden extraer las siguientes conclusiones:

\begin{itemize}
  \item La selección de HDBSCAN como algoritmo basado en densidad para BHAVRP se fundamenta en su capacidad para manejar clústeres de densidades variables y su robustez ante el ruido, superando las limitaciones de DBSCAN.
  \item La adaptación propuesta para el MDVRP, basada en una fase de asignación heurística por cercanía con restricciones de capacidad y reasignación individual de ruido, permite integrar el algoritmo en el contexto logístico sin comprometer la lógica de agrupamiento por densidad.
  \item La integración de HDBSCAN en la arquitectura de BHAVRP respeta los patrones de diseño existentes como \textit{Template Method} y los principios de diseño de software, garantizando la mantenibilidad y extensibilidad de la biblioteca mediante el uso de clases auxiliares especializadas.
\end{itemize}
%==============================================================
%                    CAPITULO 3
%==============================================================
\chapter{An\'alisis experimental}

En este cap\'itulo se presentan las pruebas funcionales y unitarias para las funcionalidades principales del algoritmo implementado, adem\'as de un estudio comparativo entre los resultados obtenidos por las heurísticas de construcción de BHCVRP usando los algoritmos de agrupamiento de la biblioteca y el nuevo algoritmo propuesto. Tambi\'en se realiza un an\'alisis de correlaci\'on entre las m\'etricas internas de BHAVRP y la calidad de las rutas obtenidas, a modo de validaci\'on complementaria.

Todos los experimentos se llevan a cabo en un entorno controlado, utilizando un ordenador con las siguientes especificaciones: procesador AMD Ryzen 5 5500U with Radeon Graphics de 2.10 GHz, 8 GB de memoria RAM y sistema operativo Windows 11 Home de 64 bits versión 25H2, usando el IDE IntelliJ IDEA con Java en su versi\'on 23.

\section{Pruebas funcionales}

Con el prop\'osito de garantizar la correcta integraci\'on y funcionamiento del nuevo algoritmo implementado en BHAVRP, se desarrolla un conjunto de pruebas unitarias usando el componente JUnit para validar las funciones esenciales del algoritmo HDBSCAN. Se prueban los m\'etodos \textit{distance} y \textit{getNeighbours} heredados de la clase AbstractDensity y el método propio del algoritmo \textit{findNearestDepotIndice}. Estas pruebas funcionales se realizan previamente al an\'alisis de calidad y rendimiento del algoritmo sobre instancias.

\subsection{Pruebas del m\'etodo \textit{distance} heredado de AbstractDensity}

El m\'etodo \textit{distance} de la clase AbstractDensity es fundamental para HDBSCAN, ya que calcula la distancia euclidiana entre dos puntos geogr\'aficos representados por objetos de la clase Location. Este m\'etodo se utiliza tanto en el c\'alculo de las distancias al n\'ucleo como en la construcci\'on del \'arbol de expansi\'on m\'inima. Para asegurar su correcto funcionamiento se dise\~nan pruebas unitarias que verifican:

\begin{itemize}
  \item C\'alculo correcto de distancias: se comparan distancias conocidas, como la distancia entre (0,0) y (3,4) que debe ser 5.0 por la m\'etrica euclidiana.
  \item Distancia a s\'i mismo: se verifica que la distancia de un punto a s\'i mismo sea cero, lo cual es una propiedad b\'asica de cualquier m\'etrica.
  \item Manejo de coordenadas negativas: se prueba con puntos que incluyen coordenadas negativas para asegurar que el c\'alculo no se ve afectado por el signo.
  \item Simetr\'ia: se comprueba que la distancia entre A y B sea igual a la distancia entre B y A, propiedad necesaria para cualquier funci\'on de distancia.
  \item Consistencia del tipo de distancia: se verifica que el tipo de distancia por defecto sea la euclidiana.
\end{itemize}

En la Figura~\ref{fig:prueba_distance_hdbscan} se puede apreciar el resultado proporcionado por el IDE a esta prueba.

\figura{fig19_prueba_distance_hdbscan.png}{0.85\textwidth}{Resultados de las pruebas del m\'etodo distance para HDBSCAN.}{fig:prueba_distance_hdbscan}

\subsection{Pruebas del m\'etodo \textit{getNeighbours} de la clase AbstractDensity}

El m\'etodo \textit{getNeighbours} es fundamental para el funcionamiento de HDBSCAN, ya que identifica los vecinos de un cliente dentro de un radio \(\varepsilon\) epsilon. Esta operaci\'on se utiliza durante el c\'alculo de las distancias al n\'ucleo y en la fase de expansi\'on de cl\'usteres. El m\'etodo recibe un cliente de entrada, la lista completa de clientes y el radio de vecindad, y retorna una lista con todos los clientes cuya distancia al cliente de entrada es menor o igual a \(\varepsilon\). Las pruebas unitarias dise\~nadas verifican:

\begin{itemize}
  \item Vecindad b\'asica: se comprueba que para un cliente con vecinos dentro del radio, el m\'etodo retorne exactamente los clientes esperados.
  \item Exclusi\'on del propio cliente: se verifica que el cliente de entrada no se incluya a s\'i mismo en la lista de vecinos, salvo que la implementaci\'on lo contemple expl\'icitamente.
  \item Vecindad vac\'ia: se prueba con un cliente aislado y un radio peque\~no, asegurando que la lista retornada est\'e vac\'ia pero no sea nula.
  \item L\'imite del radio: se comprueba que un cliente cuya distancia es exactamente igual a \(\varepsilon\) se incluya como vecino con condici\'on \(\leq\).
  \item Radio cero: se verifica que con \(\varepsilon = 0\), solo se incluyan los clientes con coordenadas id\'enticas al cliente de entrada.
  \item Manejo de coordenadas negativas: se prueba con clientes en el segundo y tercer cuadrante del plano cartesiano.
\end{itemize}

En la Figura~\ref{fig:prueba_getNeighbours} se puede apreciar el resultado proporcionado por el IDE a esta prueba.

\figura{fig23_prueba_getNeighbours.png}{0.85\textwidth}{Resultados de las pruebas del m\'etodo getNeighbours.}{fig:prueba_getNeighbours}

\subsection{Pruebas del m\'etodo \textit{findNearestDepotIndice} de la clase HDBSCAN}

El m\'etodo \textit{findNearestDepotIndice} construye una cola de prioridad con los \'indices de los dep\'ositos ordenados de menor a mayor distancia respecto a un cliente dado. Esta operaci\'on es clave en la fase de reasignaci\'on individual, ya que permite seleccionar eficientemente el dep\'osito m\'as cercano con capacidad disponible. Las pruebas unitarias dise\~nadas verifican:

\begin{itemize}
  \item Ordenamiento correcto: se comprueba que al extraer los elementos de la cola, los \'indices correspondan a los dep\'ositos en orden ascendente de distancia al cliente.
  \item Tama\~no de la cola: se verifica que la cola contenga exactamente tantos elementos como dep\'ositos haya en la instancia.
  \item Dep\'osito m\'as cercano: se comprueba que el primer elemento extra\'ido corresponda al \'indice del dep\'osito m\'as cercano al cliente.
  \item Dep\'osito m\'as lejano: se verifica que el \'ultimo elemento de la cola sea el \'indice del dep\'osito m\'as alejado.
  \item Cliente en coordenadas de un dep\'osito: se prueba con un cliente cuya ubicaci\'on coincide exactamente con un dep\'osito, verificando que ese dep\'osito sea el primero en la cola.
  \item Consistencia entre llamadas: se comprueba que dos llamadas consecutivas con el mismo cliente retornen colas con el mismo orden.
  \item Manejo de coordenadas negativas: se prueba con clientes y dep\'ositos en cuadrantes negativos.
\end{itemize}

En la Figura~\ref{fig:prueba_findNearestDepot} se puede apreciar el resultado proporcionado por el IDE a esta prueba.

\figura{fig24_prueba_findNearestDepot.png}{0.85\textwidth}{Resultados de las pruebas del m\'etodo findNearestDepotIndice.}{fig:prueba_findNearestDepot}

\section{Pruebas unitarias de caja negra}

Adem\'as de las pruebas funcionales de los m\'etodos esenciales, se realizaron pruebas de caja negra para verificar el comportamiento global del algoritmo HDBSCAN en escenarios extremos, asegurando que la adaptaci\'on al MDVRP respeta las restricciones de capacidad y asignaci\'on.

En la Tabla~\ref{tab:pruebas_hdbscan}, HDBSCAN demuestra que su mecanismo de expansi\'on controlada desde los dep\'ositos, junto con la reasignaci\'on final de clientes no asignados, logra respetar los l\'imites de capacidad incluso cuando la demanda de los cl\'usteres supera la capacidad disponible. La capacidad de HDBSCAN para identificar puntos de ruido permite tratar de forma diferenciada a los clientes at\'ipicos, mejorando la calidad estructural de las asignaciones. Estos resultados confirman que la adaptaci\'on cumple con las reglas de negocio del MDVRP.

\clearpage
\begin{landscape}
\begin{table}[htbp]
\centering
\caption{Pruebas unitarias de caja negra del algoritmo HDBSCAN.}
\label{tab:pruebas_hdbscan}
\footnotesize
\renewcommand{\arraystretch}{1.2}
\begin{tabular}{@{}cp{2.8cm}p{2.2cm}p{3.2cm}p{5cm}p{5cm}@{}}
\toprule
\textbf{No} & \textbf{Descripci\'on} & \textbf{Caso prueba} & \textbf{Juego de Datos} & \textbf{Resultados esperados} & \textbf{Resultados Reales} \\
\midrule
1 & Asignaci\'on b\'asica con dos grupos bien separados. & Asignaci\'on b\'asica & 2 dep\'ositos, 6 clientes (3 por dep\'osito). & Cl\'usteres [1,2,3] y [4,5,6]. Sin no asignados. & Cl\'usteres [1,2,3] y [4,5,6]. Sin no asignados. \\
2 & Detecci\'on y reasignaci\'on de ruido. & Detecci\'on de ruido & 7 clientes, uno aislado en (50,50). & Grupos [1,2,3] y [4,5,6]. Cliente 7 reasignado. & Grupos [1,2,3] y [4,5,6]. Cliente 7 reasignado. \\
3 & Respeto a la capacidad de los dep\'ositos. & Capacidad limitada & 2 dep\'ositos cap. 15, 6 clientes demanda 10 (total 60 $>$ 30). & Se asignan hasta 30 unidades. Resto no asignados. & Se asignan hasta 30 unidades. Resto no asignados. \\
4 & Un solo grupo denso de clientes. & Un solo cl\'uster & 5 clientes cerca del dep\'osito 1; dep\'osito 2 lejos. & 1 cl\'uster con 5 clientes. Dep\'osito 2 sin asignaci\'on. & 1 cl\'uster con 5 clientes. Dep\'osito 2 sin asignaci\'on. \\
5 & Sin estructura de densidad clara. & Todos at\'ipicos & 5 clientes dispersos sin agrupaciones. & Finaliza sin errores. Asignaci\'on m\'inima o nula. & Finaliza sin errores. Asignaci\'on m\'inima o nula. \\
6 & Configuraci\'on din\'amica de \texttt{minPoints}. & C\'alculo de \texttt{minPoints} & 6 clientes. & \texttt{minPoints}=3, \texttt{minClusterSize}=2. & \texttt{minPoints}=3, \texttt{minClusterSize}=2. \\
7 & Excepci\'on sin problema cargado. & Error en los Datos & Sin instancia cargada. & Lanza \texttt{AssignmentException}. & Lanza \texttt{AssignmentException}. \\
\bottomrule
\end{tabular}
\end{table}
\end{landscape}
\clearpage

\section{Comparaci\'on entre los algoritmos de agrupamiento}

Con el objetivo de analizar el comportamiento y la calidad de las asignaciones generadas por los algoritmos implementados en BHAVRP, se dise\~na un experimento enfocado en instancias utilizadas en trabajos anteriores \parencite{ericramosaragonIncorporacionModuloMetricas2025, aragonBHAVRPBibliotecaHeuristicas2025}. Este experimento busca comparar el rendimiento de cada algoritmo de asignaci\'on en condiciones controladas, evaluando los resultados mediante el conjunto de m\'etricas internas incluidas dentro de la biblioteca. El estudio experimental se lleva a cabo empleando las 33 instancias originalmente introducidas por \parencite{Cordeau2012} en su trabajo sobre el MDVRP. Estas instancias tienen amplio reconocimiento en el \'ambito cient\'ifico de la optimizaci\'on de rutas. La Tabla~\ref{tab:instancias_hdbscan} detalla, para cada instancia del problema, el n\'umero total de clientes, la cantidad de dep\'ositos, la flota de veh\'iculos asignada a cada dep\'osito y su capacidad.

Con el objetivo de fortalecer el an\'alisis experimental, se incorporan 11 instancias adicionales, alcanzando un total de 44 instancias para la evaluaci\'on. Estas instancias provienen de la investigaci\'on realizada en [63], y se detallan en la Tabla~\ref{tab:nuevas_instancias}.

Para el experimento se usaron todoa los algoritmos de BHAVRP junto con el nuevo algoritmo HDBSCAN, excluyendo RandomSequentialCyclic y SequentialCyclic ya que presentan problemas con algunas instancias al generar un solo clúster. Cada algoritmo se ejecuta una sola vez por instancia, utilizando como tipo de distancia aproximada la distancia euclidiana.

\begin{table}[H]
\centering
\caption{Descripci\'on de las instancias.}
\label{tab:instancias_hdbscan}
\small
\begin{tabular}{@{}ccccc@{}}
\toprule
\textbf{ID} & \textbf{\shortstack{Total de\\clientes}} & \textbf{\shortstack{Total de\\dep\'ositos}} & \textbf{\shortstack{Total de veh\'iculos\\por dep\'osito}} & \textbf{\shortstack{Capacidad\\de los veh\'iculos}} \\
\midrule
p01 & 50 & 4 & 4 & 80 \\
p02 & 50 & 4 & 2 & 160 \\
p03 & 75 & 5 & 3 & 140 \\
p04 & 100 & 2 & 8 & 100 \\
p05 & 100 & 2 & 5 & 200 \\
p06 & 100 & 3 & 6 & 100 \\
p07 & 100 & 4 & 4 & 100 \\
p08 & 249 & 2 & 14 & 500 \\
p09 & 249 & 3 & 12 & 500 \\
p10 & 249 & 4 & 8 & 500 \\
p11 & 249 & 5 & 6 & 500 \\
p12, p13, p14 & 80 & 2 & 5 & 60 \\
p15, p16, p17 & 160 & 4 & 5 & 60 \\
p18, p19, p20 & 240 & 6 & 5 & 60 \\
p21, p22, p23 & 360 & 9 & 5 & 60 \\
pr1 & 48 & 4 & 1 & 200 \\
pr2 & 96 & 4 & 2 & 195 \\
pr3 & 144 & 4 & 3 & 190 \\
pr4 & 192 & 4 & 4 & 185 \\
pr5 & 240 & 4 & 5 & 180 \\
pr6 & 288 & 4 & 6 & 175 \\
pr7 & 72 & 1 & 6 & 200 \\
pr8 & 144 & 2 & 6 & 190 \\
pr9 & 216 & 6 & 3 & 180 \\
pr10 & 288 & 6 & 4 & 170 \\
\bottomrule
\end{tabular}
\end{table}

\begin{table}[H]
\centering
\caption{Descripci\'on de las nuevas instancias.}
\label{tab:nuevas_instancias}
\small
\begin{tabular}{@{}ccccc@{}}
\toprule
\textbf{ID} & \textbf{\shortstack{Total de\\clientes}} & \textbf{\shortstack{Total de\\dep\'ositos}} & \textbf{\shortstack{Total de veh\'iculos\\por dep\'osito}} & \textbf{\shortstack{Capacidad de los\\veh\'iculos}} \\
\midrule
gj1 & 50 & 2 & 2 & 492 \\
gj2 & 75 & 3 & 3 & 415 \\
gj3 & 100 & 2 & 2 & 841 \\
gj4 & 150 & 4 & 4 & 651 \\
gj5 & 199 & 5 & 5 & 758 \\
gj6 & 75 & 2 & 2 & 669 \\
gj7 & 125 & 3 & 3 & 768 \\
gj8 & 249 & 4 & 4 & 1219 \\
gj9 & 60 & 2 & 2 & 579 \\
gj10 & 80 & 3 & 3 & 507 \\
gj11 & 120 & 5 & 5 & 1000 \\
\bottomrule
\end{tabular}
\end{table}

%Para analizar la calidad de las asignaciones generadas por cada algoritmo, se utilizan todas las m\'etricas de validaci\'on interna incorporadas en BHAVRP. Estas m\'etricas permiten evaluar cuantitativamente aspectos clave como la cohesi\'on y la separaci\'on entre cl\'usteres. En la Tabla~\ref{tab:metricas_hdbscan}, se presentan las m\'etricas a utilizar con sus siglas, el criterio que eval\'uan y una descripci\'on.

\subsection{An\'alisis estad\'istico}

A partir de los experimentos realizados, se decide utilizar la prueba estad\'istica no param\'etrica de Friedman N $\times$ N \parencite{friedmanUseRanksAvoid1937}, la cual posibilita enfrentar por pares todos los algoritmos examinados.

Para la prueba se utilizan las seis heurísticas de construcción implementadas en BHCVRP excluyendo RandomByElement debido a que presentaba errores con varias instancias. Algunas de estas heurísticas requieren de parámetros en su ejecución. A continuación, se presenta las heurística de construcción utilizadas y la configuración de estos parámetro:
\begin{itemize}
 \item Heurística del Vecino más Cercano con Lista de Candidatos Restringidos (NN): propone la construcción de una solución insertando al cliente más cercano dentro de la lista de candidatos restringida al último cliente ya insertado. Posee como parámetro el tamaño de la lista de candidatos restringidos (size\_RLC): 3
 \item Algoritmo de Ahorros en su versión secuencial (SaveSeq): este algoritmo plantea que en una solución dos rutas diferentes pueden ser combinadas formando una nueva ruta. Esta versión del algoritmo propone la construcción de las rutas una a una. Utiliza el parámetro $\lambda$ (parameter\_shape) con valor 1 para la construcción de la matriz de ahorros. 
 \item Algoritmo de Ahorros en su versión paralela (SaveParall): esta otra versión del algoritmo tiene la misma filosofía, pero trabaja sobre todas las rutas de forma simultánea. Utiliza el parámetro $\lambda$ (parameter\_shape) con valor 1 para la construcción de la matriz de ahorros.
 \item Heurística de Inserción Secuencial de Mole \& Jameson (MJ): esta heurística de inserción utiliza dos medidas para decidir el próximo cliente a insertar en la solución. Utiliza los parámetros $\lambda$ (c1) con valor 1 que representa el peso del costo directo entre los clientes en la ruta, $\mu$ (c2) con valor 1 que representa el peso del costo desde el depósito hasta el cliente que se está considerando insertar y la forma de seleccionar el primer cliente a insertar en la ruta (select\_customer\_MJ\_type) con valor 1 que significa el más cercano al depósito.
 \item Heurística de Barrido (Sweep): esta heurística trabaja en dos fases, primero forma los grupos de clientes y luego construyen las rutas para cada grupo.
 \item Heurística de Inserción en Paralelo de Christofides, Mingozzi y Toth (CMT): esta heurística de inserción opera en dos fases. En la primera fase se determina la cantidad de rutas a utilizar, junto con un cliente para inicializar cada ruta. En la segunda fase se crean dichas rutas y se inserta el resto de los clientes en ellas según el costo de inserción. Utiliza el parámetro $\lambda$ (parameter\_L) con valor 1, que controla el peso del costo de viajar desde el cliente j hasta el cliente i en el cálculo del costo de inserción.
\end{itemize}


En la Tabla~\ref{tab:friedman_hdbscan} se presenta, para cada algoritmo de asignaci\'on, su valor promedio por cada heur\'istica de construcci\'on en todas las instancias. Las celdas en verde reflejan las mejores posiciones, mientras que las celdas en rojo indican los desempe\~nos m\'as bajos.

\clearpage
	\begin{landscape}
\begin{table}[H]
\centering
\caption{Resultados de la prueba de Friedman NxN.}
\label{tab:friedman_hdbscan}
\footnotesize
\begin{tabular}{@{}lcccccccccc@{}}
\toprule
\textbf{Algoritmo} & \textbf{CMT} & \textbf{MoleJameson} & \textbf{NearestNeighborWithRLC} & \textbf{SaveParallel} & \textbf{SaveSequential} & \textbf{Sweep} \\
\midrule
    BestCyclicAssignment & 12.2727 & 12.375 & 9.9545 & 12.2955 & 12.4205 & 10.3636 \\
    BestNearest & 8.5909 & 6.7841 & 6.25  & 8.0341 & 7.8523 & 8.7727 \\
    CLARA & 11.7727 & 11    & 12.1364 & 11.9545 & 11.9091 & 12.3409 \\
    CoefficientPropagation & 11.9318 & 12.4318 & 10.1364 & 10.7045 & 12.0455 & 12.1591 \\
    CURE  & \cellcolor[rgb]{ 0,  .69,  .314}6.0682 & \cellcolor[rgb]{ 0,  .69,  .314}5.7841 & 4.5682 & 6.4091 & 6.5568 & 8.2045 \\
    CyclicAssignment & 16.2045 & 15.6591 & 14.6591 & 15.8409 & 15.7273 & 15.1591 \\
    DBSCAN & 11.8409 & 12.3182 & 10.7045 & 10.9545 & 12.4318 & 9.2386 \\
    Farthest-First & 10.1136 & 9.4659 & 7.2727 & 10.8409 & 9.8864 & 7.9091 \\
    HDBSCAN & 8.1818 & 8.7955 & 9.5   & 7.5795 & 8.7955 & 10.5227 \\
    KMEANS & 6.1818 & 8.0795 & 5.2727 & 7.6477 & 8.1818 & 7.4545 \\
    NearestByCustomer & 6.8636 & 6.6023 & 13.0455 & 7.8409 & 7.1818 & 6.8068 \\
    NearestByDepot & 13.0455 & 11.9545 & 9.3636 & 13.3977 & 12.4432 & 11.3636 \\
    PAM   & 11.5682 & 12.4773 & 10.2045 & 11.6364 & 10.9773 & 12.6591 \\
    Parallel & 7.25  & 6.8523 & 12.7273 & 7.8409 & 7.7159 & \cellcolor[rgb]{ 0,  .69,  .314}6.2727 \\
    RandomByElement & \cellcolor[rgb]{ 1,  0,  0}19 & \cellcolor[rgb]{ 1,  0,  0}18.9773 & \cellcolor[rgb]{ 1,  0,  0}19 & \cellcolor[rgb]{ 1,  0,  0}19 & \cellcolor[rgb]{ 1,  0,  0}19 & \cellcolor[rgb]{ 1,  0,  0}18.9318 \\
    Simplified & 7.2727 & 7.6364 & 10.3182 & 7.75  & 6.6818 & 8.4091 \\
    Sweep & 7.2045 & 6.1477 & 11.0682 & 6.6818 & 6.2614 & 6.5227 \\
    ThreeCriteriaClustering & 7.9545 & 7.8295 & 9.4659 & 9.1136 & 7.875 & 8.8409 \\
    UPGMC & 6.6818 & 8.8295 & \cellcolor[rgb]{ 0,  .69,  .314}4.3523 & \cellcolor[rgb]{ 0,  .69,  .314}4.4773 & \cellcolor[rgb]{ 0,  .69,  .314}6.0568 & 8.0682 \\
\bottomrule
\end{tabular}
\end{table}
	\end{landscape}
\clearpage

A continuaci\'on, se presenta la interpretaci\'on de los resultados obtenidos en el ranking:

\begin{itemize}
  \item Entre los mejores algoritmos se posicionan CURE, UPGMC, KMEANS y Sweep lideran el ranking en la mayor\'ia de las heur\'isticas de construcci\'on. CURE obtiene el mejor resultado en heurística como CMT y MoleJameson, mientras que UPGMC es el mejor en NearestNeighborWithRLC, SaveParallel y SaveSequential, mientras que KMEANS y Sweep mantienen posiciones competitivas de forma consistente.
  
  \item HDBSCAN obtiene un ranking promedio que lo ubica en una posici\'on intermedia dentro del conjunto de algoritmos evaluados. Sus mejores resultados se dan con SaveParallel y CMT, mientras que su peor desempe\~no ocurre con Sweep. Aunque no alcanza los primeros puestos, HDBSCAN supera consistentemente a DBSCAN en todas las heur\'isticas y muchas de las heurísticas clásicas de asignación, lo cual confirma la mejora que aporta respecto a su predecesor dentro de la misma familia de algoritmos basados en densidad.
  
  \item RandomByElement obtiene el peor ranking de forma sistem\'atica, seguido por CyclicAssignment y NearestByDepot. CLARA, BestCyclicAssignment, PAM, CoefficientPropagation y DBSCAN tambi\'en ocupan posiciones desfavorables.
\end{itemize}

Es importante destacar que HDBSCAN, aunque no sea de los más destacados, presenta un comportamiento s\'olido y consistente en los resultados de las heurísticas de construcción, con un desempe\~no significativamente superior al de DBSCAN y al de los algoritmos con peores resultados. Su principal fortaleza radica en la detecci\'on y manejo del ruido.

Una vez obtenidos los rankings promedio de Friedman, se aplicaron los procedimientos \textit{post-hoc} de Bergmann \parencite{bergmannImprovementsGeneralMultiple1988}, Holm \parencite{holmSimpleSequentiallyRejective1979} y Shaffer \parencite{shafferModifiedSequentiallyRejective1986} para controlar la tasa de error familiar. La prueba de Friedman arroj\'o un estad\'istico $\chi^2 = 313.15$ con 18 grados de libertad y un valor $p = 1.55 \times 10^{-10}$, lo que confirma la existencia de diferencias estad\'isticamente significativas entre los algoritmos.

Los resultados del an\'alisis \textit{post-hoc} permiten extraer las siguientes conclusiones sobre el algoritmo HDBSCAN:

\begin{itemize}
  \item Es significativamente mejor que RandomByElement con ($p = 0$), lo cual confirma que las asignaciones generadas por HDBSCAN producen rutas de mejor calidad que las del algoritmo aleatorio.
  \item Es significativamente mejor que DBSCAN con ($p = 0.0049$), demostrando que la evoluci\'on jer\'arquica basada en estabilidad supera al DBSCAN cl\'asico en el contexto de la asignaci\'on del MDVRP.
  \item Es significativamente mejor que NearestByDepot con ($p = 0.000001$), superando a la heur\'istica de asignaci\'on por dep\'osito m\'as cercano.
  \item Es significativamente mejor que PAM y CLARA con ($p = 0.00072$ y $p = 0.00027$ respectivamente), superando a los algoritmos particionales cl\'asicos.
  \item UPGMC es significativamente mejor que HDBSCAN ($p = 0.0097$), lo cual es consistente con el ranking promedio.
  \item No se encontraron diferencias estad\'isticamente significativas con CURE, KMEANS, BestNearest y Sweep, lo que indica que HDBSCAN es competitivo con los mejores algoritmos en t\'erminos de calidad de rutas, aunque su promedio es ligeramente inferior.
\end{itemize}

En s\'intesis, HDBSCAN demuestra ser un algoritmo que aporta mejoras significativas respecto a DBSCAN y a las heur\'isticas cl\'asicas, aunque no supera a los mejores algoritmos del estado de la biblioteca como CURE y UPGMC en el ranking global. Su principal valor reside en la calidad estructural de las asignaciones y en su capacidad para manejar el ruido, lo cual no siempre se traduce directamente en mejores rutas.

\section{An\'alisis de correlaci\'on entre m\'etricas internas y calidad de rutas}

Con el objetivo de evaluar si existe una relaci\'on entre la calidad estructural de las asignaciones, medida mediante las m\'etricas internas de BHAVRP que se pueden apreciar en la Tabla~\ref{tab:metricas}, y la calidad de las rutas obtenidas en la fase posterior, se realiza un an\'alisis de correlaci\'on complementario. A diferencia del estudio de \parencite{ericramosaragonIncorporacionModuloMetricas2025}, donde no se encontr\'o correlaci\'on significativa entre la suma de las métricas normalizadas, en este estudio se analiza la relaci\'on para cada métrica y heur\'istica de construcci\'on por separado. 

\clearpage
\begin{landscape}
\begin{table}[htbp]
\centering
\caption{Resumen de las m\'etricas incorporadas en BHAVRP.}
\label{tab:metricas}
\small
\begin{tabular}{@{}lp{1.3cm}p{2cm}p{8.5cm}@{}}
\toprule
\textbf{M\'etrica} & \textbf{Siglas} & \textbf{?`Qu\'e eval\'ua?} & \textbf{Descripci\'on} \\
\midrule
\textbf{Suma de Errores al Cuadrado} & SSE & Cohesi\'on & Suma de las distancias cuadr\'aticas entre cada punto y el centroide de su grupo. \textbf{Minimizar}. \\
\textbf{\'Indice de Ball-Hall} & BH & Cohesi\'on & Mide la dispersi\'on promedio dentro de los cl\'usteres. \textbf{Minimizar}. \\
\textbf{Suma de Cuadrados entre Cl\'usteres} & SSB & Separaci\'on & Eval\'ua la separaci\'on entre los grupos considerando la distancia entre sus centroides y el centro global de los datos. \textbf{Maximizar}. \\
\textbf{\'Indice de Calinski-Harabasz} & CH & Cohesi\'on y Separaci\'on & Relaci\'on entre la dispersi\'on entre cl\'usteres (SSB) y la dispersi\'on dentro de los cl\'usteres (SSE). \textbf{Maximizar}. \\
\textbf{\'Indice de Dunn} & DI & Cohesi\'on y Separaci\'on & Relaci\'on entre la distancia m\'inima inter-cl\'uster y el di\'ametro m\'aximo intra-cl\'uster (el tama\~no del grupo m\'as grande). \textbf{Maximizar}. \\
\textbf{\'Indice de Davies-Bouldin} & DB & Cohesi\'on y Separaci\'on & Analiza la relaci\'on entre la dispersi\'on interna de los cl\'usteres y la distancia entre sus centroides. \textbf{Minimizar}. \\
\textbf{Coeficiente de Silueta} & CS & Cohesi\'on y Separaci\'on & Combina cohesi\'on y separaci\'on para cada punto. \textbf{Maximizar}. \\
\textbf{\'Indice de Xie-Beni} & XB & Cohesi\'on y Separaci\'on & Penaliza los agrupamientos donde los puntos est\'an alejados de sus centroides y los que est\'an muy pr\'oximos entre s\'i. \textbf{Minimizar}. \\
\textbf{\'Indice de Ray-Turi} & RT & Cohesi\'on y Separaci\'on & Eval\'ua la relaci\'on entre la dispersi\'on dentro de los cl\'usteres y la distancia m\'inima entre los centroides de los cl\'usteres. \textbf{Minimizar}. \\
\bottomrule
\end{tabular}
\end{table}
\end{landscape}
\clearpage

Para cada heur\'istica de construcci\'on se calcularon los coeficientes de Spearman y Pearson entre cada m\'etrica interna y el costo total de las rutas. Los resultados completos se presentan en las Tablas~\ref{tab:correlacion_saveparallel} a \ref{tab:correlacion_nnrlc1}.

\begin{table}[H]
\centering
\caption{Correlaci\'on entre m\'etricas internas y costo total (heur\'istica SaveParallel).}
\label{tab:correlacion_saveparallel}
\small
\begin{tabular}{@{}lcccc@{}}
\toprule
\textbf{M\'etrica} & \textbf{Spearman} & \textbf{p-valor} & \textbf{Pearson} & \textbf{p-valor} \\
\midrule
BH & 0.8555 & $3.02\times10^{-6}$ & 0.9161 & $3.68\times10^{-8}$ \\
CH & -0.6921 & $1.02\times10^{-3}$ & -0.7436 & $2.63\times10^{-4}$ \\
DB & 0.7577 & $1.72\times10^{-4}$ & 0.8459 & $5.06\times10^{-6}$ \\
DI & -0.7893 & $5.88\times10^{-5}$ & -0.8474 & $4.68\times10^{-6}$ \\
RT & 0.6816 & $1.31\times10^{-3}$ & 0.7578 & $1.71\times10^{-4}$ \\
CS & -0.6605 & $2.08\times10^{-3}$ & -0.7123 & $6.23\times10^{-4}$ \\
SSB & -0.6611 & $2.06\times10^{-3}$ & -0.8084 & $2.81\times10^{-5}$ \\
SSE & 0.6945 & $9.69\times10^{-4}$ & 0.7575 & $1.73\times10^{-4}$ \\
XB & 0.6804 & $1.35\times10^{-3}$ & 0.7541 & $1.92\times10^{-4}$ \\
\bottomrule
\end{tabular}
\end{table}

\begin{table}[H]
\centering
\caption{Correlaci\'on entre m\'etricas internas y costo total (heur\'istica SaveSequential).}
\label{tab:correlacion_savesequential}
\small
\begin{tabular}{@{}lcccc@{}}
\toprule
\textbf{M\'etrica} & \textbf{Spearman} & \textbf{p-valor} & \textbf{Pearson} & \textbf{p-valor} \\
\midrule
BH & 0.8648 & $1.78\times10^{-6}$ & 0.9520 & $3.63\times10^{-10}$ \\
CH & -0.7305 & $3.83\times10^{-4}$ & -0.8337 & $9.24\times10^{-6}$ \\
DB & 0.7749 & $9.77\times10^{-5}$ & 0.8945 & $2.39\times10^{-7}$ \\
DI & -0.7942 & $4.90\times10^{-5}$ & -0.8335 & $9.33\times10^{-6}$ \\
RT & 0.7217 & $4.86\times10^{-4}$ & 0.8371 & $7.85\times10^{-6}$ \\
CS & -0.7199 & $5.10\times10^{-4}$ & -0.8015 & $3.70\times10^{-5}$ \\
SSB & -0.6292 & $3.90\times10^{-3}$ & -0.8533 & $3.42\times10^{-6}$ \\
SSE & 0.7380 & $3.09\times10^{-4}$ & 0.8451 & $5.26\times10^{-6}$ \\
XB & 0.7222 & $4.79\times10^{-4}$ & 0.8348 & $8.79\times10^{-6}$ \\
\bottomrule
\end{tabular}
\end{table}

\begin{table}[H]
\centering
\caption{Correlaci\'on entre m\'etricas internas y costo total (heur\'istica Sweep).}
\label{tab:correlacion_sweep}
\small
\begin{tabular}{@{}lcccc@{}}
\toprule
\textbf{M\'etrica} & \textbf{Spearman} & \textbf{p-valor} & \textbf{Pearson} & \textbf{p-valor} \\
\midrule
BH & 0.8068 & $2.99\times10^{-5}$ & 0.8612 & $2.20\times10^{-6}$ \\
CH & -0.7717 & $1.09\times10^{-4}$ & -0.8159 & $2.06\times10^{-5}$ \\
DB & 0.7723 & $1.07\times10^{-4}$ & 0.8513 & $3.80\times10^{-6}$ \\
DI & -0.7126 & $6.18\times10^{-4}$ & -0.7728 & $1.05\times10^{-4}$ \\
RT & 0.7682 & $1.22\times10^{-4}$ & 0.8253 & $1.36\times10^{-5}$ \\
CS & -0.7366 & $3.22\times10^{-4}$ & -0.7716 & $1.09\times10^{-4}$ \\
SSB & -0.6617 & $2.03\times10^{-3}$ & -0.8165 & $2.00\times10^{-5}$ \\
SSE & 0.7775 & $8.93\times10^{-5}$ & 0.8176 & $1.91\times10^{-5}$ \\
XB & 0.7723 & $1.07\times10^{-4}$ & 0.8237 & $1.47\times10^{-5}$ \\
\bottomrule
\end{tabular}
\end{table}

\begin{table}[H]
\centering
\caption{Correlaci\'on entre m\'etricas internas y costo total (heur\'istica CMT).}
\label{tab:correlacion_cmt}
\small
\begin{tabular}{@{}lcccc@{}}
\toprule
\textbf{M\'etrica} & \textbf{Spearman} & \textbf{p-valor} & \textbf{Pearson} & \textbf{p-valor} \\
\midrule
BH & 0.9315 & $6.92\times10^{-9}$ & 0.9522 & $3.47\times10^{-10}$ \\
CH & -0.8235 & $1.48\times10^{-5}$ & -0.8540 & $3.28\times10^{-6}$ \\
DB & 0.8495 & $4.19\times10^{-6}$ & 0.9156 & $3.86\times10^{-8}$ \\
DI & -0.8126 & $2.36\times10^{-5}$ & -0.8383 & $7.41\times10^{-6}$ \\
RT & 0.8042 & $3.32\times10^{-5}$ & 0.8600 & $2.35\times10^{-6}$ \\
CS & -0.7981 & $4.23\times10^{-5}$ & -0.8226 & $1.54\times10^{-5}$ \\
SSB & -0.7301 & $3.86\times10^{-4}$ & -0.8829 & $5.59\times10^{-7}$ \\
SSE & 0.8284 & $1.18\times10^{-5}$ & 0.8624 & $2.05\times10^{-6}$ \\
XB & 0.8004 & $3.87\times10^{-5}$ & 0.8580 & $2.63\times10^{-6}$ \\
\bottomrule
\end{tabular}
\end{table}

\begin{table}[H]
\centering
\caption{Correlaci\'on entre m\'etricas internas y costo total (heur\'istica MoleJameson).}
\label{tab:correlacion_molejameson}
\small
\begin{tabular}{@{}lcccc@{}}
\toprule
\textbf{M\'etrica} & \textbf{Spearman} & \textbf{p-valor} & \textbf{Pearson} & \textbf{p-valor} \\
\midrule
BH & 0.7612 & $1.54\times10^{-4}$ & 0.9245 & $1.54\times10^{-8}$ \\
CH & -0.7717 & $1.09\times10^{-4}$ & -0.8569 & $2.81\times10^{-6}$ \\
DB & 0.7301 & $3.86\times10^{-4}$ & 0.8864 & $4.37\times10^{-7}$ \\
DI & -0.6547 & $2.35\times10^{-3}$ & -0.7569 & $1.76\times10^{-4}$ \\
RT & 0.7630 & $1.45\times10^{-4}$ & 0.8548 & $3.14\times10^{-6}$ \\
CS & -0.7507 & $2.13\times10^{-4}$ & -0.8218 & $1.60\times10^{-5}$ \\
SSB & -0.7021 & $8.06\times10^{-4}$ & -0.8662 & $1.64\times10^{-6}$ \\
SSE & 0.7705 & $1.13\times10^{-4}$ & 0.8592 & $2.46\times10^{-6}$ \\
XB & 0.7617 & $1.51\times10^{-4}$ & 0.8547 & $3.16\times10^{-6}$ \\
\bottomrule
\end{tabular}
\end{table}

\begin{table}[H]
\centering
\caption{Correlaci\'on entre m\'etricas internas y costo total (heur\'istica NearestNeighborWithRLC).}
\label{tab:correlacion_nnrlc}
\small
\begin{tabular}{@{}lcccc@{}}
\toprule
\textbf{M\'etrica} & \textbf{Spearman} & \textbf{p-valor} & \textbf{Pearson} & \textbf{p-valor} \\
\midrule
BH & 0.5426 & $1.64\times10^{-2}$ & 0.6879 & $1.13\times10^{-3}$ \\
CH & -0.4583 & $4.85\times10^{-2}$ & -0.6088 & $5.67\times10^{-3}$ \\
DB & 0.5274 & $2.03\times10^{-2}$ & 0.6908 & $1.06\times10^{-3}$ \\
DI & -0.6959 & $9.36\times10^{-4}$ & -0.8299 & $1.11\times10^{-5}$ \\
RT & 0.4241 & $7.04\times10^{-2}$ & 0.6254 & $4.19\times10^{-3}$ \\
CS & -0.4943 & $3.15\times10^{-2}$ & -0.6249 & $4.23\times10^{-3}$ \\
SSB & -0.3975 & $9.19\times10^{-2}$ & -0.6177 & $4.83\times10^{-3}$ \\
SSE & 0.4642 & $4.52\times10^{-2}$ & 0.6236 & $4.33\times10^{-3}$ \\
XB & 0.4133 & $7.86\times10^{-2}$ & 0.6200 & $4.63\times10^{-3}$ \\
\bottomrule
\end{tabular}
\end{table}

\begin{table}[H]
\centering
\caption{Correlaci\'on entre m\'etricas internas y costo total (heur\'istica NearestNeighborWithRLC con RLC=1).}
\label{tab:correlacion_nnrlc1}
\small
\begin{tabular}{@{}lcccc@{}}
\toprule
\textbf{M\'etrica} & \textbf{Spearman} & \textbf{p-valor} & \textbf{Pearson} & \textbf{p-valor} \\
\midrule
BH & 0.2828 & $2.41\times10^{-1}$ & 0.2447 & $3.13\times10^{-1}$ \\
CH & -0.1353 & $5.81\times10^{-1}$ & -0.1692 & $4.89\times10^{-1}$ \\
DB & 0.2555 & $2.91\times10^{-1}$ & 0.2648 & $2.73\times10^{-1}$ \\
DI & -0.5610 & $1.25\times10^{-2}$ & -0.6209 & $4.56\times10^{-3}$ \\
RT & 0.1036 & $6.73\times10^{-1}$ & 0.1863 & $4.45\times10^{-1}$ \\
CS & -0.1827 & $4.54\times10^{-1}$ & -0.1870 & $4.43\times10^{-1}$ \\
SSB & -0.0061 & $9.80\times10^{-1}$ & -0.1465 & $5.50\times10^{-1}$ \\
SSE & 0.1484 & $5.44\times10^{-1}$ & 0.1885 & $4.40\times10^{-1}$ \\
XB & 0.0957 & $6.97\times10^{-1}$ & 0.1805 & $4.60\times10^{-1}$ \\
\bottomrule
\end{tabular}
\end{table}

Los resultados revelan que todas las m\'etricas internas presentan una correlaci\'on estad\'isticamente significativa con el costo total de las rutas con (valor $p < 0.05$) en la mayor\'ia de las heur\'isticas de construcci\'on evaluadas. La \'unica excepci\'on es la heur\'istica NearestNeighborWithRLC con RLC=1, donde solo el \'Indice de Dunn (DI) mantiene significaci\'on estad\'istica. Este resultado contrasta con la conclusi\'on del estudio previo, donde se afirmaba que no exist\'ia correlaci\'on. 

\subsection*{Interpretaci\'on de las correlaciones}

Las correlaciones observadas siguen la direcci\'on esperada seg\'un el criterio de cada m\'etrica:

\begin{itemize}
  \item M\'etricas de cohesi\'on (SSE, BH): presentan correlaci\'on {positiva} con el costo. Es decir, valores m\'as altos de dispersi\'on intra-cl\'uster se asocian con rutas m\'as costosas. El \'Indice de Ball-Hall (BH) destaca con la correlaci\'on m\'as fuerte (Pearson hasta $0.952$ en SaveSequential y CMT), lo que lo posiciona como el mejor predictor del costo dentro de este conjunto.
  \item M\'etricas de separaci\'on (SSB): presentan correlaci\'on {negativa}. Una mayor separaci\'on entre cl\'usteres se asocia con un menor costo de rutas.
  \item M\'etricas combinadas (CH, DI, CS): muestran correlaci\'on {negativa} con el costo. Los \'indices de Dunn (DI) y Calinski-Harabasz (CH) destacan por su capacidad predictiva (DI: Pearson hasta $-0.847$ en SaveParallel; CH: Pearson hasta $-0.857$ en MoleJameson).
  \item \'Indices DB, RT y XB: presentan correlaci\'on {positiva}. Valores m\'as altos se asocian con costos m\'as elevados.
\end{itemize}

\subsection*{Patrones por heur\'istica}

Al comparar las heur\'isticas de construcci\'on, se observan patrones consistentes:

\begin{itemize}
  \item {SaveSequential, CMT y MoleJameson} presentan las correlaciones m\'as fuertes, con m\'etricas como BH y DB alcanzando coeficientes superiores a $0.9$.
  \item {Sweep y SaveParallel} muestran correlaciones s\'olidas y significativas en todas las m\'etricas.
  \item {NearestNeighborWithRLC} es la heur\'istica con correlaciones m\'as d\'ebiles, y algunas de ellas no resultan significativas (RT, SSB, XB con Spearman $p > 0.05$).
  \item {NearestNeighborWithRLC=1} presenta el caso m\'as extremo: solo el \'Indice de Dunn mantiene correlaci\'on significativa. Esto sugiere que con un tama\~no de lista de candidatos restringidos muy peque\~no, las asignaciones pierden estructura y las m\'etricas internas dejan de ser predictivas del costo.
\end{itemize}

En todos los casos, {BH (Ball-Hall), DI (Dunn), CH (Calinski-Harabasz), DB (Davies-Bouldin) y SSE} destacan como las m\'etricas con mayor poder predictivo del costo total. Esto las posiciona como candidatas ideales para ser utilizadas como criterios de selecci\'on de algoritmos de asignaci\'on cuando no se dispone del costo de las rutas.

\section{Conclusiones parciales}

A partir del estudio experimental realizado en este cap\'itulo, se pueden extraer las siguientes conclusiones:

\begin{itemize}
  \item Las pruebas funcionales verificaron el correcto funcionamiento de los m\'etodos esenciales de HDBSCAN, incluyendo el c\'alculo de distancias, la obtención de vecinos y la búsqueda del depósito más cercano, garantizando su integraci\'on fiable en BHAVRP.
  
  \item Las pruebas unitarias de caja negra confirmaron que la adaptaci\'on de HDBSCAN respeta las restricciones de capacidad y maneja adecuadamente escenarios extremos como datos err\'oneos, dep\'ositos vac\'ios y cl\'usteres no asignables.
  
  \item La comparaci\'on estad\'istica mediante la prueba de Friedman y los an\'alisis \textit{post-hoc} revel\'o que CURE, UPGMC, KMEANS y Sweep conforman el grupo de mejor desempe\~no en la calidad de las rutas. HDBSCAN se posiciona en un nivel intermedio, superando significativamente a DBSCAN y a las heur\'isticas cl\'asicas, y sin diferencias significativas con varios de los mejores algoritmos, lo que demuestra su rendimiento y su aporte como alternativa basada en densidad.
  
  \item El an\'alisis de correlaci\'on demostr\'o que las m\'etricas internas de validaci\'on presentan correlaci\'on estad\'isticamente significativa con el costo total de las rutas en la mayor\'ia de las heur\'isticas de construcci\'on. Los \'indices de Ball-Hall y Dunn destacan como los mejores predictores, lo que confirma la utilidad pr\'actica de las m\'etricas internas para seleccionar algoritmos de asignaci\'on sin necesidad de construir las rutas.
\end{itemize}

\printbibliography[title={Referencias Bibliográficas}, heading=bibintoc]

\end{document}